\pdfoutput=1
\documentclass[11pt]{amsart}
\usepackage{amsmath,amssymb,amsthm,mathtools,mathrsfs}
\usepackage[margin=1.2in]{geometry}
\usepackage{booktabs,array,longtable,enumitem,xcolor}
\usepackage[colorlinks=true,linkcolor=blue,citecolor=blue,urlcolor=blue]{hyperref}
\hypersetup{pdfauthor={Bernd Johannes Wuebben},
 pdftitle={Cobordism maps and Atiyah–Floer isomorphisms for framed instanton homology II},
 pdfsubject={A conditional integral cobordism comparison with transported orientations},
 pdfkeywords={Atiyah-Floer conjecture, framed instanton homology, determinant lines, coherent orientations, cobordism maps}}
\setlength{\emergencystretch}{2em}
\newtheorem{theorem}{Theorem}[section]
\newtheorem{proposition}[theorem]{Proposition}
\newtheorem{lemma}[theorem]{Lemma}
\newtheorem{corollary}[theorem]{Corollary}
\theoremstyle{definition}
\newtheorem{definition}[theorem]{Definition}
\newtheorem{hypothesis}[theorem]{Hypothesis}
\newtheorem{convention}[theorem]{Convention}
\newtheorem{example}[theorem]{Example}
\theoremstyle{remark}
\newtheorem{remark}[theorem]{Remark}
\theoremstyle{plain}
\newtheorem{thmletter}{Theorem}
\renewcommand{\thethmletter}{\Alph{thmletter}}
\numberwithin{equation}{section}
\newcommand{\ZZ}{\mathbb Z}
\newcommand{\RR}{\mathbb R}
\newcommand{\CC}{\mathbb C}
\newcommand{\FF}{\mathbb F_2}
\newcommand{\cB}{\mathcal B}
\newcommand{\cC}{\mathcal C}
\newcommand{\cD}{\mathcal D}
\newcommand{\cE}{\mathcal E}
\newcommand{\cF}{\mathcal F}
\newcommand{\cG}{\mathcal G}
\newcommand{\cH}{\mathcal H}
\newcommand{\cJ}{\mathcal J}
\newcommand{\cK}{\mathcal K}
\newcommand{\cL}{\mathcal L}
\newcommand{\cM}{\mathcal M}
\newcommand{\cP}{\mathcal P}
\newcommand{\cU}{\mathcal U}
\newcommand{\cW}{\mathcal W}
\newcommand{\bB}{\mathbf B}
\newcommand{\SO}{\mathrm{SO}}
\newcommand{\SU}{\mathrm{SU}}
\newcommand{\U}{\mathrm U}
\newcommand{\CF}{\mathit{CF}}
\newcommand{\HF}{\mathit{HF}}
\newcommand{\Ish}{I^{\#}}
\newcommand{\IG}{I^{\#}_{J}}
\newcommand{\Msf}{\mathsf M}
\newcommand{\id}{\mathrm{id}}
\newcommand{\eps}{\varepsilon}
\newcommand{\ad}{\operatorname{ad}}
\newcommand{\ind}{\operatorname{ind}}
\newcommand{\coker}{\operatorname{coker}}
\newcommand{\im}{\operatorname{im}}
\newcommand{\Hom}{\operatorname{Hom}}
\newcommand{\End}{\operatorname{End}}
\newcommand{\Aut}{\operatorname{Aut}}
\newcommand{\tr}{\operatorname{tr}}
\newcommand{\ev}{\operatorname{ev}}
\newcommand{\ten}{\operatorname{ten}}
\newcommand{\PD}{\operatorname{PD}}
\newcommand{\cont}{\operatorname{cont}}
\newcommand{\orZ}{\mathfrak o}
\title[Atiyah--Floer cobordism maps II]{Cobordism maps and Atiyah--Floer isomorphisms\\
for framed instanton homology II}
\author{Bernd Johannes Wuebben}
\date{26 September 2026}
\subjclass[2020]{57K31, 57R58 (primary); 53D40, 58J05 (secondary)}
\keywords{Atiyah--Floer conjecture, framed instanton homology, coherent
orientations, determinant lines, Lagrangian Floer homology, cobordism maps}
\begin{document}
\begin{abstract}
We study the extension to integer coefficients of the cobordism
comparison between framed instanton homology and the Lagrangian Floer
homology of a Heegaard splitting. Orientations are transported from
almost-complex orientations on the gauge side through moduli spaces of
the mixed anti-self-duality and holomorphic curve equation. Subject to
explicit analytic hypotheses on completed mixed domains and oriented
gluing, we obtain a projective natural isomorphism for connected
cobordisms with a framed arc and a bundle trivialized at the ends.
The symplectic maps are signed counts of the triangles, compression
cups and caps used in the mod-two comparison. The principal additional
arguments concern the constancy of sphere-gluing signs, reflection of
coherent orientations, and the identification of the diagonal
compression cup with continuation. We keep the grading signs in the
gauge orientation convention explicit. The result allows one overall
sign per cobordism; comparison with homology-oriented integral instanton
maps and with relative-spin quilt orientations is left open.
\end{abstract}
\maketitle
\begin{center}
\small\textit{First draft. The principal comparison is conditional on
Hypotheses~\ref{hyp:regularity}, \ref{hyp:completion}, and
\ref{hyp:selfgluing}. Their verification remains part of the integral
extension.}
\end{center}
\setcounter{tocdepth}{1}
\tableofcontents

\section{Introduction}\label{sec:introduction}

\subsection{The integral comparison}
Let $M$ be a closed connected oriented three-manifold. The framed
instanton construction replaces $M$ by $M^\#=M\#T^3$ and uses an
$\SO(3)$ bundle which is trivial on $M$ and odd on a torus in the
added summand. A Heegaard splitting determines two Lagrangians in the
moduli space of flat connections on the stabilized surface. Daemi,
Fukaya and Lipyanskiy construct an isomorphism between the resulting
gauge and symplectic Floer homologies by counting solutions of a mixed
equation~\cite{DFL}. The mixed equation couples the anti-self-duality
equation on a four-dimensional region to a holomorphic curve equation
on a surface; the two solutions have matching flat boundary values.

Paper~I~\cite{PaperI} concerns compatibility with four-dimensional
cobordisms over $\FF$. Its symplectic maps count holomorphic triangles
and compression cups and caps. The two analytic comparisons there are
a collapse of a compression body next to a holomorphic strip and an
excision along a torus in a mixed domain. The purpose of this paper is
to give the orientation arguments required for integer coefficients.
We retain those geometric maps: the symplectic cobordism map is not
defined by conjugating an instanton map through an isomorphism.

There are three distinct orientation questions. First, the quotient
by the lifting class used in framed instanton homology must preserve
orientations. Second, the collapse and torus-excision identifications
must preserve the orientations of individual solutions. Third, the
orientation system transported through a reflected mixed domain must
be compared with the dual of the original system. The last comparison
is needed to turn the tensor supplied by the compression cup into a
map with the source complex used in Paper~I. An unspecified sign on
each generator would not suffice: the ambiguity in a cobordism map
must be a single sign.

We use the almost-complex orientation construction of
Kronheimer--Mrowka~\cite[Section~5.1]{KM-unknot}. It is invariant under
bundle automorphisms. The almost-complex structure on a cobordism is
fixed near the odd torus, but its restrictions at the two ends are
allowed to vary. Conversion operators then identify the end
orientation lines with fixed object lines. Their placement relative
to the generator lines matters when a conversion is odd. We denote
the resulting gauge functor by $\IG$ to record this convention.

\subsection{Statement and hypotheses}
The category $\bB$ is described in Section~\ref{sec:geometry}. An object
is a closed connected oriented three-manifold with the framing data
needed for its torus summand. A morphism is a connected oriented
cobordism $W$ equipped with a normally framed arc $\nu$ joining its
ends and an $\SO(3)$ bundle $E$ trivialized at the ends and along the
arc. A presentation $p$ consists of a Heegaard splitting, metric data,
small primary holonomy perturbations, and regularizing data. The
additional smallness and reflection conventions for integer
coefficients are specified in Section~\ref{sec:choices}.

Write $C_G(M,p)$ for the integral gauge complex with the almost-complex
orientation convention, and $\CF(M,p)$ for the symplectic complex with
transported orientations. The graded homotopy category $\cK(\ZZ)$
and the notation $F\simeq\pm G$ are defined in
Section~\ref{sec:algebra}. For an admissible decomposition $d$ of a
cobordism, $S_d(W)$ is the composite of signed holomorphic triangle
maps, compression cups, compression caps, and continuation maps.

The geometric arguments use the unsigned compactness and gluing
results of Paper~I and the following additional hypotheses. They are
stated precisely in Section~\ref{sec:analytic}.
\begin{enumerate}[label=(\roman*)]
\item Uniform regularity and compatible Fredholm realizations on the
completed mixed domains, including families, changes of weights, and
the constant reference configurations.
\item An oriented torus-completion comparison, compatible with
ordinary gluing and with the reference configurations used to
normalize the vertex maps.
\item Determinant-one self-gluing along the nonseparating hypersurface
in the gauge excision cobordism, with its gluing multiplicities and
determinant-line identifications.
\end{enumerate}
These are mathematical hypotheses, not assertions that follow merely
from the existence of the corresponding mod-two count. In particular,
the first draft does not claim an unconditional integral version of
the mixed excision theorem.

\begin{thmletter}[Conditional integral comparison]\label{thm:A}
Assume the analytic results of Paper~I and
Hypotheses~\ref{hyp:regularity}, \ref{hyp:completion}, and
\ref{hyp:selfgluing}. With the orientation and admissibility
conventions of this paper, the mixed maps
\[
 N_p:C_G(M,p)\longrightarrow\CF(M,p)
\]
are isomorphisms of finite free integral complexes. For every morphism
$(W,E,\nu):M_0\to M_1$, admissible endpoint presentations $p_0,p_1$,
and admissible decomposition $d$, one has
\begin{equation}\label{eq:main}
 N_{p_1}\IG(W,E,\nu)\simeq
 \pm S_d(W,E,\nu)N_{p_0}
 \quad\text{in }\cK(\ZZ).
\end{equation}
The sign is one sign for the entire map. Consequently the
symplectic maps are independent of the decomposition up to sign and
chain homotopy. On homology they define a projective functor and
$H(N)$ is a projective natural isomorphism from $\IG$ to that functor.
Each symplectic homology map has the parity of the corresponding
gauge map.

At objects with nonzero integral gauge homology the symplectic maps
are the geometric maps of Paper~I with the orientations constructed
here. If the integral gauge homology vanishes, the complex is
contractible; maps incident to it represent the zero morphism in
$\cK(\ZZ)$, with the convention of
Section~\ref{subsec:zero}. Reduction modulo two recovers the maps of
Paper~I, up to chain homotopy in the contractible case.
\end{thmletter}

The qualifier ``projective'' has the usual literal meaning: maps are
identified with their negatives. No system of signs making all
composition identities exact is asserted. In particular, multiplying
individual elementary maps by signs is permitted, whereas multiplying
different matrix entries by unrelated signs is not.

\subsection{The two orientation comparisons}
The proof of Theorem~\ref{thm:A} separates two further statements.
They use the same hypotheses as that theorem.

\begin{thmletter}[Orientations under collapse]\label{thm:B}
The bijection of regular index-zero solutions in the compression-body
collapse of Paper~I preserves the transported determinant-line
orientations. This assertion uses the metric and the admissible
linearizations specified in Section~\ref{sec:determinants}. It is
compatible with ordinary gauge and strip gluing and with compact
parameter families. The two-handle comparison therefore holds over
$\ZZ$, up to the single sign arising from the auxiliary orientation
data of the cobordism.
\end{thmletter}

\begin{thmletter}[Reflection and the diagonal cup]\label{thm:C}
Fix a Heegaard splitting, metric data, and an almost-complex base
structure at an object with nonzero integral gauge homology. In the
admissible ball of primary perturbations, the sphere character of
the transported strip orientation system is constant and agrees
with that of the reflected presentation. There is consequently a
reflection isomorphism
\[
 \Phi_p:\CF(M,\bar p)\longrightarrow\CF(M,p)^\vee
\]
determined by the configuration-level orientations up to one sign.
With this pairing, the compression cup with no handle attached is
chain homotopic, up to one sign, to the symplectic continuation map.
\end{thmletter}

The sphere character records the change of a coherent orientation
when a sphere is attached to a strip configuration. Its constancy
does not compute its value. The reflection statement compares two
characters; it does not assert that either character is trivial.
This distinction avoids imposing an unnecessary relative-spin
orientation convention on the compression correspondence.

\subsection{Strategy}
The integral mixed map has a unit diagonal in the filtration used
in Paper~I. The diagonal configuration is a constant pair, and its
orientation is the canonical orientation of an invertible operator.
The resulting integral isomorphism is stronger than a mod-two
isomorphism: an odd integer coefficient on the diagonal would not
be enough.

For two-handles, determinant lines are transported through the
collapse family. The relevant continuity is continuity over smooth
configurations with all derivatives controlled on compact subsets
and with the prescribed asymptotic behavior. Strong convergence of
solutions alone does not identify all operator families used in an
orientation argument. We state the needed realization and completion
properties explicitly and prove the sign comparison from them.

For one-handles, a folded triangle gives a tensor in the product of
the reflected source complex and the target complex. Excision
identifies this tensor with a gauge handle state. This gives a cup
formula before any reflection pairing is inserted. The remaining
problem is to identify the pairing with the one for the actual
symplectic continuation maps. We first prove that the sphere
character is constant under small changes of primary data, then
compare it across a fold. A rigidity argument for coherent systems
identifies the diagonal cup with continuation. The cup formula now
gives small-primary naturality, mixed duality, and the one-handle
square. Reflection gives the three-handle square.

Finally, the algebraic descent uses chain homotopies throughout.
Over $\ZZ$, equality on homology is generally insufficient to prove
that two chain maps are homotopic. The proof therefore retains the
homotopies supplied by the one-dimensional moduli spaces and uses
the invertibility of the mixed maps only after the elementary
squares have been established.

\subsection{Scope of the draft}
The almost-complex convention used for $\IG$ is part of the statement.
We do not identify it with Scaduto's homology-oriented integral
cobordism maps~\cite{Scaduto}. An odd conversion can introduce a
grading operator rather than an overall sign. Section~\ref{sec:gauge}
explains this point and the additional grading information needed
for a comparison. We likewise do not identify $N_p$ with the integral
map of~\cite{DFL}; their reductions modulo two are compared in
Paper~I at common admissible data.

The orientations on holomorphic curves are transported from the
gauge side. Comparison with the relative-spin orientations used in
quilted Floer theory~\cite{WW-orientations,WW-quilt} is a separate
question. The underlying curves and correspondences are the same
ones as in Paper~I, but an equality of mod-two counts does not supply
this integral orientation comparison. Disconnected objects, units,
counits, closed pairings, and nontrivial bundles on the objects are
outside the present statement. Bundles on cobordisms may be
nontrivial relative to the specified trivialized ends.

\section{Geometric data and analytic hypotheses}\label{sec:geometry}

\subsection{Framed objects and cobordisms}
Fix an oriented ball in $M$ and form $M^\#=M\#T^3$ along that ball.
Write $T^\circ=T^3\setminus\operatorname{int}B^3$. Choose a torus
$T_0\subset T^\circ$, a neighborhood $N(T_0)$, and an odd bundle
$P_T\to T^3$ with $\langle w_2(P_T),[T_0]\rangle=1$. The bundle
$P_M\to M^\#$ is obtained by joining $P_T$ to the trivial bundle on
$M$. The lifting class
\[
 \alpha_M=\PD[T_0]\pmod2\in H^1(M^\#;\ZZ/2)
\]
acts on the determinant-one configuration space; the framed complex
uses the quotient by $\phi_M=\langle\alpha_M\rangle$.

For $(W,E,\nu)$ as in the introduction, remove a tube around the
normally framed arc and insert $T^\circ\times[0,1]$. Denote the
result by $(W^\#,E^\#)$. The lifting subgroup $\phi_W$ is the one
induced by this torus tube, with the endpoint restrictions $\phi_{M_i}$.
At an intermediate disconnected cut we retain the full subgroup
specified by the tube; it need not be the product of the groups on
the components. This distinction is essential in the excision
calculation in Section~\ref{sec:excision}.

The category identifies isomorphic decorated cobordisms and uses
gluing along the fixed endpoint data for composition. The connected
source and target condition permits a handle decomposition with
one-, two-, and three-handles, together with changes of presentation.
As in Paper~I, elementary words also include metric changes,
Heegaard moves and frame changes. The presentation graph is
connected and the gauge cobordism map is independent of such a
word. These are the geometric inputs to the descent at the end of
the paper.

\subsection{Compression bodies and the symplectic complex}
If $M=H_0\cup_\Sigma H_1$ is a Heegaard splitting of genus $g$,
adjoin $[0,1]\times T^2$ to each handlebody. The resulting torus
compression bodies $C_i$ have boundary $\Sigma_h\sqcup T^2$, where
$h=g+1$. Let $\Msf_h$ be the moduli space of flat connections on the
odd $\SO(3)$ bundle over $\Sigma_h$, with the lift convention of
\cite[Section~2]{DFL}. Its symplectic form is the usual
Atiyah--Bott form. Restriction of flat connections on $C_i$ gives a
Lagrangian $L(C_i)\subset\Msf_h$.

Primary perturbations $h_i$ are finite sums of averaged holonomy
functions supported in $C_i^\circ$, away from the fixed collars.
Their critical equations are
\begin{equation}\label{eq:pertflat}
 *_3F_B+\nabla_Bh_i=0.
\end{equation}
We use a sign-sensitive trace convention for fixed-determinant
$\U(2)$ holonomy: on each solid torus choose a continuous square
root of the determinant holonomy and take the trace of the resulting
$\SU(2)$ part. The choices are made once for the finite family.
Equation~\eqref{eq:pertflat} determines perturbed Lagrangians
$P_i=L_{h_i}(C_i)$. At transverse data, their intersection points
correspond to the gauge critical orbits.

For each intersection point $x$, let $a_x$ be the corresponding
gauge critical point and $L_G(a_x)$ its real determinant line.
For a real line $L$, write $\orZ(L)$ for the rank-one free abelian
group generated by its two orientations with the relation
$[-o]=-[o]$. The integral symplectic chain group is
\begin{equation}\label{eq:chains}
 \CF(M,p)=\bigoplus_{x\in P_0\cap P_1}\orZ(L_G(a_x)).
\end{equation}
The differential counts rigid strips with the transported
orientation system of Section~\ref{sec:vertex}. The mod-two
reduction of~\eqref{eq:chains} is the chain group of Paper~I.

\subsection{Mixed configurations}
A mixed domain consists of a gauge region and a symplectic surface
meeting along a matching boundary. On the gauge side a connection
satisfies a perturbed anti-self-duality equation. On the symplectic
side a map to $\Msf_h$, or to a product of such spaces, satisfies a
perturbed Cauchy--Riemann equation. Along a matching boundary the
surface restriction of the connection is flat and represents the
value of the map. At ordinary ends we prescribe irreducible
nondegenerate gauge limits or transverse Lagrangian intersection
points. At mixed ends we retain the finite-dimensional stationary
coordinates explicitly.

The domains used below are the mixed vertex domain, mixed and folded
triangles, their strip gluings, the closed-cut mixed domain, the
finite-width torus domains, their completed limits, and the collapse
families. We also use ordinary strips, continuation strips, and the
disks obtained by doubling and capping a folded triangle. Statements
about determinant lines apply on configuration spaces, not just on
their solution subsets.

We use the smooth topology on these spaces: in fixed end
trivializations all derivatives converge on compact sets, the
asymptotic limits converge with their parameters, and the decaying
parts converge in the chosen weighted norms at every derivative
order. A continuous compact family means a family continuous in
this topology. Sobolev completions are used locally to construct
Fredholm operators and determinant lines. They do not replace the
smooth topology when a holonomy family is differentiated repeatedly.

\subsection{The analytic assumptions}\label{sec:analytic}
The unsigned moduli spaces, their energy-index estimates, and the
low-dimensional compactifications are those of Paper~I. The
following assumptions specify the additional uniform and oriented
statements used in this draft. The ordinary local Fredholm and
regularity results for the mixed equation are those of
\cite{DFL,DFL-mixed}; the hypotheses below include their use on the
additional completed domains of this paper.

\begin{hypothesis}[Uniform realizations and regularity]\label{hyp:regularity}
For each mixed domain listed above, and for each compact parameter
family used below, the following hold.
\begin{enumerate}[label=(\alph*)]
\item Smooth configurations with fixed asymptotics admit local
matching-boundary charts. The gauge-fixed linearization on weighted
spaces is Fredholm; its kernels and cokernels are smooth. Changes
of chart and the admissible changes of realization of
Definition~\ref{def:realization} give continuous families of
Fredholm operators and the same determinant line.
\item At ordinary ends the weights lie in common spectral gaps.
At mixed ends the stationary coordinates are separated from the
decaying part. Regular zero-dimensional solutions and the constant
reference pairs have invertible realized operators. On compact
families the estimates, slice radii, and smallness thresholds can
be chosen uniformly, including at infinite torus length.
\item The mixed shifting construction, its relative family form,
and node gluing are continuous in the smooth topology. They commute
with ordinary end gluing. A compact family of shifted configurations
can be made exactly symplectically constant outside one common
compact set, without changing a member already constant there.
\item The solution and homotopy counts are independent of the
choice of weights within the specified gaps. The pure symplectic
continuation, cap, and strip-gluing moduli spaces have the usual
regular low-dimensional compactifications, with no additional
bubbling boundary in the dimensions used here.
\end{enumerate}
\end{hypothesis}

The relevant asymptotic restriction is strict: the weight is less
than the smallest positive decay rate of every end occurring in
that instance. Uniformity is asserted for a fixed compact family,
not over all Heegaard splittings or all primary perturbations.
When an end is completed, the same statement must hold on the
completed domain. A finite-neck estimate by itself does not imply
this last assertion.

\begin{hypothesis}[Oriented mixed torus completion]\label{hyp:completion}
Consider the torus domains in the proof of the mixed excision
theorem of Paper~I, with fixed admissible primary and secondary
data. The determinant lines extend continuously across the
width and completion comparisons. Their identifications have the
following properties:
\begin{enumerate}[label=(\alph*)]
\item they are natural under the lifting groups and compatible with
gauge and strip gluing, including the conversion lines at the
closed cut;
\item at a split constant reference pair they are the product of
the two canonical invertible-operator orientations;
\item the preglued configurations and the solutions are connected
by paths in the stated charts which identify these determinant
lines, with the same property for the one-dimensional families
used to compare lengths and data;
\item the compactified one-dimensional comparison has only the
broken ends and endpoint strata used in Paper~I. With outward
normal first and the parameter orientation fixed, its boundary
gives the integral chain homotopy asserted in
Theorem~\ref{thm:mixed-excision}.
\end{enumerate}
\end{hypothesis}

This assumption concerns an oriented comparison of a particular
family of mixed equations. It does not assume the cup formula,
reflection invariance, or the cobordism comparison. It is precisely
the extension of the torus-completion argument for which a complete
integral verification is still needed.

\begin{hypothesis}[Determinant-one self-gluing]\label{hyp:selfgluing}
Let $X$ be a gauge cobordism occurring in the torus-excision
construction, and cut it along its designated closed hypersurface
$Y$. Suppose the limiting critical points on $Y$ are irreducible
and nondegenerate and the data are translation invariant on the
neck. For sufficiently large neck length the index-zero solutions
on $X$ are obtained by gluing the matched determinant-one solutions
on the cut cobordism. The fiber of this gluing has the lifting
multiplicity described in Lemma~\ref{lem:gluing-multiplicity}.
There are no additional index-zero limits. The determinant-line
gluing uses the graded symmetry of Section~\ref{sec:algebra}; its
sign depends only on the ordered end degrees and the fixed
orientation choices. The corresponding one-dimensional families
give the composition homotopies.
\end{hypothesis}

Ordinary composition and trace constructions in instanton Floer
theory provide the models for this hypothesis
\cite[Sections~5.2--5.3]{KM-unknot},
\cite[Section~2.5]{KM-sutures}. Here its nonseparating form and
the precise determinant-one multiplicity are part of the assumption.
We do not infer this form from a trace in a quotient complex.

\section{Complexes, graded signs and determinant lines}\label{sec:algebra}

\subsection{The homotopy category}
A complex is a finite free $\ZZ/2$-graded abelian group
$C=C^0\oplus C^1$ with an odd differential $d_C$. For a homogeneous
map $F:C\to D$ of degree $f$, set
\begin{equation}\label{eq:hom-differential}
 \delta F=d_DF-(-1)^fFd_C.
\end{equation}
A closed map satisfies $\delta F=0$. Thus an odd closed map
anticommutes with the differentials. We allow sums of the two
degrees when discussing homotopies of the composites below.
The category $\cK(\ZZ)$ has these complexes as objects and closed
maps modulo boundaries as morphisms. In particular
\[
 F\simeq G\quad\Longleftrightarrow\quad F-G=\delta K.
\]
For homogeneous $F,G$ of degree $f$, the homotopy $K$ can be
chosen of degree $f+1$. The notation $F\simeq\pm G$ means that
one of $F-G,F+G$ is a boundary, with a single choice of sign.

If $A$ and $B$ are closed and $B$ has degree $b$, then
\begin{equation}\label{eq:compose-homotopy}
 (\delta K)A=\delta(KA),\qquad
 B\delta K=(-1)^b\delta(BK).
\end{equation}
These identities justify composition of every homotopy used below.
They also show that projectivizing by the relation $[F]=[-F]$
is compatible with composition.

\begin{lemma}\label{lem:acyclic}
An acyclic finite free $\ZZ/2$-graded complex over $\ZZ$ is
contractible. An even quasi-isomorphism between finite free
complexes is a chain homotopy equivalence.
\end{lemma}
\begin{proof}
Let $B^i=\im(d:C^{i-1}\to C^i)$. Each $B^i$ is free as a
subgroup of a finitely generated free group. If the complex is
acyclic, the exact sequence
$0\to B^i\to C^i\xrightarrow{d}B^{i+1}\to0$ splits. Relative
to these splittings the differential identifies the complementary
summand of $C^i$ with $B^{i+1}$. Its inverse on this summand,
and zero on the complement, defines $h$ with $dh+hd=1$.
Apply this to the mapping cone of a quasi-isomorphism to obtain
the second assertion.
\end{proof}

The lemma does not say that a map inducing zero on homology is
null-homotopic. Nor does it identify arbitrary free complexes
with their homology, which may contain torsion. We use the lemma
only for acyclic complexes and actual quasi-isomorphisms.

\subsection{The tensor convention}
Generator lines precede the determinant lines of trajectories
and cobordisms. With this ordering the tensor differential is
\begin{equation}\label{eq:tensor}
 d(x\otimes y)=(-1)^{|y|}dx\otimes y+x\otimes dy.
\end{equation}
For homogeneous maps $F,G$ we use
\begin{equation}\label{eq:tensor-maps}
 (F\otimes_R G)(x\otimes y)
   =(-1)^{|F||y|}Fx\otimes Gy.
\end{equation}
Closed maps tensor to a closed map of degree $|F|+|G|$ with
this rule. The usual left tensor differential is related to
\eqref{eq:tensor} by the even chain isomorphism
\[
 T(x\otimes y)=(-1)^{|x||y|}x\otimes y.
\]
Thus a formula imported in a left tensor convention is conjugated
by $T$ at each tensor end before it is used here.

\begin{lemma}\label{lem:tensor-geometric}
For a disconnected gauge end $Y_1\sqcup Y_2$, with the generator
line ordered as $L(a_1)\otimes L(a_2)$, gluing index-one
trajectories on the two components gives~\eqref{eq:tensor}.
\end{lemma}
\begin{proof}
A trajectory on $Y_1$ has an odd determinant line. To place it
next to $L(a_1)$, move it past $L(a_2)$; this gives
$(-1)^{|a_2|}$. A trajectory on $Y_2$ is already adjacent to
its generator line. Summing the two contributions gives the
formula.
\end{proof}

Put $\eps_Cx=(-1)^{|x|}x$. If a homogeneous map $F$ of degree
$f$ commutes strictly with the differentials, the map
$F^{\mathrm r}=F\eps_C^f$ is closed. For strictly commuting
maps $F,G$ of degrees $f,g$, respectively,
\begin{equation}\label{eq:right-conversion}
 (FG)^{\mathrm r}=(-1)^{fg}F^{\mathrm r}G^{\mathrm r}.
\end{equation}
The discrepancy is one sign per composite. By contrast the
factor $\eps_C$ is not generally an overall sign.

For example, on $C^0=C^1=\ZZ$ with zero differential, the map
interchanging the two generators and its product with $\eps_C$
are not equal up to sign. The difference persists on homology.
This elementary example explains why an integral comparison
must specify the convention for odd cobordism maps.

\subsection{Duals, evaluation and tensors of maps}
The graded dual has
\[
 d^\vee f=-(-1)^{|f|}f\circ d,
 \qquad F^\vee(f)=(-1)^{|F||f|}f\circ F.
\]
The canonical even double-dual map is
\[
 \ev_C:C\longrightarrow C^{\vee\vee},\qquad
 \ev_C(x)(f)=(-1)^{|x||f|}f(x).
\]
The sign is necessary: the unsigned double-dual identification
does not commute with the differential in this convention.
There is an unsigned chain isomorphism
\begin{equation}\label{eq:beta}
 \beta:C^\vee\otimes_R D\longrightarrow\Hom(C,D),\qquad
 \beta(f\otimes y)(x)=f(x)y.
\end{equation}
Indeed, substituting the dual differential in~\eqref{eq:tensor}
gives exactly~\eqref{eq:hom-differential} after evaluation on a
homogeneous $x$. For a tensor in $C\otimes_R D$ we write
\[
 B_C=\beta\circ(\ev_C\otimes_R1):
 C\otimes_R D\longrightarrow\Hom(C^\vee,D).
\]
If $e_i$ is a homogeneous basis of $C$, the tensor corresponding
to a homogeneous map $F:C\to D$ is
\[
 \ten(F)=\sum_i e_i^\vee\otimes F(e_i).
\]
This expression is independent of the chosen basis. It is a
cycle if and only if $F$ is closed. For $F=1$, it is the
coevaluation tensor in the present convention.

\subsection{Determinant lines}
For a real Fredholm operator $D:X\to Y$ set
\[
 \det D=\Lambda^{\max}\ker D\otimes
       (\Lambda^{\max}\coker D)^*,\qquad
 |\det D|=\ind D\pmod2.
\]
We use the graded determinant convention of
\cite{Zinger}. Stabilization by a finite-dimensional map
$V\to Y$ with $D\oplus V$ surjective gives the local
trivializations. Enlarging the stabilization gives compatible
transition maps. This construction applies to continuous
families of Fredholm operators without requiring constant
kernel dimension.

For direct sums the determinant lines are ordered in the
order of the summands. Interchanging two summands gives
$(-1)^{\ind D_1\ind D_2}$. The dual of a tensor reverses
the order of its factors and uses nested evaluation. An
invertible operator has the canonical determinant orientation
of $\RR$; a continuous path of invertible operators preserves
this orientation.

\begin{lemma}[Ordinary gluing]\label{lem:line-gluing}
Suppose two Fredholm problems have matching ordinary ends with
invertible asymptotic operator. Gluing with a sufficiently long
neck gives a determinant-line isomorphism
\[
 \det D_1\otimes\det D_2\longrightarrow\det(D_1\#D_2)
\]
well defined up to a positive scalar. It is continuous in
compact parameter families, natural under bundle
automorphisms, and associative for disjoint necks. On
invertible operators it is the canonical identification.
\end{lemma}
\begin{proof}
Choose finite-dimensional stabilizations supported away from
the ends so that the two stabilized operators are surjective.
Patch their right inverses using cutoffs. The error tends to
zero with the neck length, so a Neumann series corrects the
patched operator to a right inverse. Projection of the patched
kernels to the kernel of the glued stabilized operator gives
an isomorphism for long necks. Its top exterior power and the
stabilization isomorphisms define the asserted map.

Enlarging stabilizations commutes with this construction.
Changing cutoffs, projections or neck length gives a path of
kernel isomorphisms, so the sign does not change. For two
disjoint necks use stabilizations supported away from both;
the two successive patching maps are homotopic to the direct
three-piece patching map. This proves associativity, with
only the graded signs caused by an actual reordering. When
both operators are invertible, no stabilization is needed
and the map is $\RR\otimes\RR\to\RR$. The gauge version is
the determinant-line construction accompanying ordinary
Floer gluing~\cite[Proposition~5.11]{Donaldson-book}; the strip
version uses the same stabilization argument for real
Cauchy--Riemann operators~\cite{Floer-Hofer}.
\end{proof}

Gluing a constant cylinder is the positive identity: its
operator is invertible, and the patched kernel map converges
to the identity as the neck grows. Weighted gluing is treated
by conjugating with the exponential weight. The resulting
asymptotic operator must be invertible, and matched ends must
use compatible weights. We will never use ordinary gluing
at a mixed end without separately keeping its stationary
coordinates.

\section{The integral gauge convention}\label{sec:gauge}

\subsection{Almost-complex orientations with free ends}
Choose an $S^1$-invariant almost-complex germ $J_T$ near
$S^1\times N(T_0)$, with $T_0$ an almost-complex torus in its
specified orientation. On an object, such a structure is equivalently
specified by a nonvanishing vector field on $M^\#$, normalized near
$T_0$. Choose a base $J^0_M$ for each object. On a cobordism require
$J_W$ to agree with the fixed germ near $N(T_0)\times[0,1]$ and to be
translation invariant on its end collars. Its end restrictions
$J'_0,J'_1$ are not prescribed in advance.

\begin{lemma}\label{lem:almost-complex}
Every decorated cobordism admits an almost-complex structure of
this kind.
\end{lemma}
\begin{proof}
The fiber of the bundle of compatible oriented almost-complex
structures is $\SO(4)/\U(2)\cong S^2$. Let
$A=N(T_0)\times[0,1]$. The primary relative obstruction lies in
$H^3(W^\#,A;\ZZ)$ and maps to $W_3(W^\#)=0$. Here an oriented
four-manifold admits a spin$^c$ structure; equivalently its second
Stiefel--Whitney class has an integral lift. The remaining relative
class is in the image of the connecting map from $H^2(A;\ZZ)$.
The restriction $H^2(W^\#;\ZZ)\to H^2(A;\ZZ)$ is surjective:
the torus summand supplies a class pairing once with $T_0$, and
this class extends by the Mayer--Vietoris sequence across the
three-sphere where the tube is attached. Thus the connecting
map vanishes.

The secondary obstruction lies in $H^4(W^\#,A;\ZZ)$.
Since $A\simeq T^2$, its third cohomology vanishes; since
$W^\#$ is connected with nonempty boundary, $H^4(W^\#;\ZZ)=0$.
The relative group therefore vanishes. Homotopy extension makes
the structure equal to the prescribed germ, and a collar
reparametrization makes it translation invariant near the ends.
\end{proof}

Prescribing both end structures would change this obstruction
problem. We use Lemma~\ref{lem:almost-complex} only with free end
restrictions. Conversion maps, rather than an extension assertion
with fixed ends, implement the fixed object convention.

The almost-complex deformation of the anti-self-duality operator
has a complex-linear endpoint. The complex orientation of its
kernel and cokernel orients its determinant line. This
construction is invariant under the full group of bundle
automorphisms, which act complex linearly at that endpoint
\cite[Section~5.1]{KM-unknot}. Ordinary determinant gluing preserves
these orientations. In particular the action of $\phi_M$ on the
orientation lines has no additional sign twist.

\subsection{The comparison character}
Let $v,v'$ be the normalized vector fields corresponding to
$J,J'$ on $S^1\times Y$, where $Y=M^\#$. Their difference class
$\delta(v,v')\in H^2(Y;\ZZ)$ satisfies
\[
 2\delta(v,v')=c_1(\xi_{v'})-c_1(\xi_v).
\]
The difference class is defined by the obstruction to homotoping
the vector fields on the two-skeleton, rather than by choosing
an arbitrary half of the displayed Chern-class difference.
This is relevant when $H^2(Y;\ZZ)$ has two-torsion.

\begin{proposition}\label{prop:character-gauge}
If $\alpha\in H^1(Y;\ZZ/2)$ is the reduction of an integral
class, the change of almost-complex orientation around a loop
with lifting class $\alpha$ is
\begin{equation}\label{eq:gauge-character}
 \xi_{J,J'}(\alpha)
   =(-1)^{\langle\delta(v,v'),\PD\alpha\rangle}.
\end{equation}
In particular it is $1$ on $\phi_M$ when $J$ and $J'$ agree
near $T_0$.
\end{proposition}
\begin{proof}
Represent the loop by a path of connections whose endpoints
are identified by a bundle automorphism of class $\alpha$.
Closing up gives a bundle $P_\alpha$ on $S^1\times Y$ with
\[
 w_2(P_\alpha)=\operatorname{pr}_Y^*w_2(P)
                  +[dt]\smile\alpha.
\]
The determinant-line comparison for this closed operator is
divided by the comparison for the constant loop. The
$P$-independent orientation factor therefore cancels.

For completeness, the bundle-dependent comparison can be
computed on a reducible model $\RR\oplus L$. The anti-self-dual
operator on the $L$-summand has both the complex structure
induced by $J$ and that of $L$. On the summand where these
structures are opposite, conjugation changes the determinant
orientation by $(-1)^{\ind_{\CC}}$. The index formula gives,
for two almost-complex structures, the ratio
\[
 (-1)^{\langle c_1(L)\smile\delta(v,v'),[S^1\times Y]\rangle}.
\]
This is the relative form of the usual almost-complex
orientation calculation\cite[Section~3(c)]{Donaldson87}.
Changing the instanton number does not change the ratio:
instanton addition preserves both complex orientations.

Because $\alpha$ has an integral lift, so does the displayed
Stiefel--Whitney class. The reducible model therefore computes
the desired ratio. Only the term $[dt]\smile\alpha$ survives
after division by the constant-loop comparison; integration
over the circle gives~\eqref{eq:gauge-character}. The
closing-up identification is made in pair form
$\det D\otimes(\det D_{\rm const})^*$, so it introduces no
choice of orientation at the closed end. It is the real
excision identification, independent of $J$ and $J'$.
Finally $\alpha_M$ is dual to $T_0$, and the difference class
restricts to zero there when the structures have the same germ.
\end{proof}

Only integral lifting classes are used in this proposition.
The lifting classes in the framed objects have that property.
No assertion for arbitrary mod-two classes is needed.

\subsection{Conversion maps and their ordering}
For a critical point $a$, let $Q_a$ interpolate the complex
operators for $J$ and $J'$. Choose a weight for which this
operator is Fredholm and choose an orientation $q_a$ of its
determinant line. Transport $q_a$ along paths of connections.
Proposition~\ref{prop:character-gauge} says that this transport
is consistent on the framed quotient. It is unique up to one
sign after a choice at one reference connection.

We place $q_a$ after the generator line and obtain
\begin{equation}\label{eq:Theta}
 \Theta^{\mathrm r}_{J\to J'}:
 C_G(Y,J)\longrightarrow C_G(Y,J'),\qquad
 x\longmapsto\gamma(x\otimes q_a).
\end{equation}
Its parity $t=\ind Q_a\pmod2$ is independent of $a$. Moving
the odd trajectory line past $q_a$ shows that
\[
 d\Theta^{\mathrm r}=(-1)^t\Theta^{\mathrm r}d.
\]
Thus it is a closed isomorphism of degree $t$. If
$\Theta^{\mathrm l}$ denotes the convention with $q_a$ placed
before the generator, then
\begin{equation}\label{eq:theta-grading}
 \Theta^{\mathrm r}=\Theta^{\mathrm l}\eps^t.
\end{equation}
Composition of conversion maps agrees with the direct
conversion up to one sign. The sign is constant because the
orientation comparison is defined over the connected space
of connections, not chosen separately at its critical points.

\begin{definition}\label{def:gauge-map}
For $(W,E,\nu)$ choose an admissible $J_W$ with end
restrictions $J'_0,J'_1$. Define
\begin{equation}\label{eq:gauge-map}
 \IG(W)=\Theta^{\mathrm r}_{J'_1\to J^0_{M_1}}
        I(W^\#,E^\#,J_W)^{\phi_W}
        \Theta^{\mathrm r}_{J^0_{M_0}\to J'_0}.
\end{equation}
At the induced end structures the middle map uses the
almost-complex orientations and the stated ordering of its
determinant line. Tensor ends use~\eqref{eq:tensor-maps}.
\end{definition}

\begin{proposition}\label{prop:gauge-functor}
The homotopy class of~\eqref{eq:gauge-map} is independent of
$J_W$ up to one sign. These classes define a projective gauge
functor, and their reduction modulo two is the framed gauge
functor of Paper~I.
\end{proposition}
\begin{proof}
Compare two choices of $J_W$ after inserting their endpoint
conversions. They orient the same real determinant line over
the space of connections with fixed endpoint limits. Their
ratio is continuous and takes values in $\{1,-1\}$, so it is
constant on that space. Gluing a path at either endpoint
shows that the ratio does not change when the limiting
critical point is changed. The character calculation makes
this assertion invariant under changing the lift of an end
critical point. Thus the ratio is one sign for all entries
of the map.

For composition choose an admissible almost-complex structure
on the composite and make it cylindrical at the intermediate
cut. The restrictions then have identical structures there,
and the ordinary gauge composition theorem applies
\cite[Proposition~5.4]{KM-unknot}. Insert the conversions to
the chosen object bases. The conversions at the common end
cancel up to sign, and~\eqref{eq:right-conversion} accounts
for their ordering. Independence of $J_W$ gives the result
for the original choices. The product cobordism gives the
identity class. All orientation signs disappear modulo two.
\end{proof}

\subsection{Reflection and the relation to other conventions}
Choose the base on the oppositely oriented object in pairs:
\[
 J^0_{-M}=\Psi_*\rho_*J^0_M.
\]
Here $\rho$ reverses the cylindrical direction and $\Psi$
is the fixed identification of the reflected torus-summed
presentation with $(-M)^\#$. Take $T_0=\{\theta_0\}\times T^2$
and the product germ there; require $\Psi$ to preserve the
odd torus, its germ, and its lifting class. Gauge reflection
then gives an even isomorphism
\begin{equation}\label{eq:gauge-duality}
 \Phi_{G,M}:C_G(-M)\longrightarrow C_G(M)^\vee,
 \qquad \Phi_{G,-M}=\Phi_{G,M}^\vee\ev.
\end{equation}
The transpose formula for a reversed cobordism follows by
bending its cylindrical ends and applying the same
ordered determinant gluing. All constant end lines are even;
any remaining change is one orientation sign of the
cobordism. This also gives the transpose formula for gauge
continuation.

If the conversion lines are instead placed on the left, put
$I_{\mathrm l}(W)=\Theta^{\mathrm l}_1I(W)\Theta^{\mathrm l}_0$.
If the two conversions have parities $t_0,t_1$, then
\begin{equation}\label{eq:gauge-dictionary}
 \IG(W)=(-1)^{t_0t_1}I_{\mathrm l}(W)\eps^{t_0+t_1}.
\end{equation}
This is an identity of maps in the stated conventions, rather
than an assertion that the two maps differ only by a sign.
The degree-dependent factor matters on an object whose
homology has both parities.

A $\ZZ/4$ grading compatible with all the orientation data
would provide a further regrading
$\eta(x)=(-1)^{g(x)(g(x)-1)/2}x$. The elementary identity
\[
 \eta(g+d)\eta(g)^{-1}
   =(-1)^{dg+d(d-1)/2}
\]
explains how an odd grading operator can then be converted
into a conjugation and one sign per arrow. Identifying such
a lift with Scaduto's data is not part of this draft. We
therefore keep $\IG$ in the theorem and make no integral
identification on the basis of the mod-two equality alone.

\section{Determinant lines on mixed configuration spaces}
\label{sec:determinants}

\subsection{Admissible operator realizations}
Let $c$ be a smooth configuration on a mixed domain, with
its specified limits. Denote the mixed equation by
$\mathcal F(c)=0$. A chart $\Phi$ centered at $c$ incorporates
the exact matching condition and the normal seam condition.
Its differential identifies a fixed linear boundary problem
with the tangent space to the configuration space.

\begin{definition}\label{def:realization}
An admissible realization of the linearized problem at $c$
is an operator of the form
\begin{equation}\label{eq:realization}
 L_c=\bigl(\mathcal S\,d\mathcal F_{\Phi(0)}\,d\Phi_0,
              \lambda^2d_{A'}^{*,g'}\bigr).
\end{equation}
It acts on the prescribed weighted spaces with stationary
mixed-end coordinates included. Here $\mathcal S$ is a
continuously chosen invertible target identification, equal
to the identity outside a compact set; $A'$ is a smooth
slice center differing from the connection in $c$ only on
a compact core and sufficiently little in $C^0$; $g'=f^2g$
with $f=1$ off the collapse region and $f$ depending only on
base coordinates near a seam; and $\lambda>0$ equals one
outside a compact set. The chart has the same end and seam
conventions as the original mixed problem. Changes of all
these choices are taken inside their fixed components.
\end{definition}

The restriction on the target identification prevents an
arbitrary reversal of orientation from being included as a
change of realization. The reference identification is the
one induced by the geometric charts. Its admissible
component is fixed before transporting orientations.

\begin{lemma}[Realization comparison]\label{lem:realization}
Under Hypothesis~\ref{hyp:regularity}, admissible realizations
at a fixed configuration have canonically identified
orientation double covers. The identification is continuous
in compact families and commutes with ordinary gluing. At a
regular zero it preserves the canonical determinant
orientation.
\end{lemma}
\begin{proof}
Changes of charts and target identifications are paths of
bounded invertible changes of source and target. They induce
the standard naturality maps on determinant lines. Varying
a weight inside a gap gives a path of Fredholm realizations;
kernel and cokernel elements have the stronger decay common
to the two weights, so their identifications are the natural
ones. For a change of slice center at a solution, split the
tangent space into infinitesimal gauge directions and a
fixed complement. Gauge invariance gives
$d\mathcal F_c\,d_{A_c}=0$. The gauge-fixed operator is then
block triangular. Its gauge block is invertible by the
Poincare estimate and the smallness of $A'-A_c$; its other
block is the intrinsic linearized equation on the quotient.
Thus invertibility at a regular zero is unchanged.

Conformal changes in the specified unperturbed region do
not change the anti-self-duality equation in four dimensions.
They do not change the surface Cauchy--Riemann equation
or the normal seam subspace. The same block argument applies
to the gauge row. Positive row multiplication does not alter
an orientation. These paths give the comparison. A
homotopy between two such paths gives an orientation cover
over a square, so its endpoint identification is unchanged.
The gluing statement follows by choosing the changes away
from the ordinary gluing collars and then using continuity.
At a regular zero every operator in these paths is
invertible, so the canonical orientation is preserved.
\end{proof}

The lemma uses the uniform analytic assertions in
Hypothesis~\ref{hyp:regularity}; it is not a proof of their
validity on a completed domain. In particular, moving a
Coulomb center at a configuration which is not a solution
is not justified by the block-triangular argument above.
We will use a two-parameter determinant-line comparison
when such a configuration occurs in the collapse argument.

\subsection{Shifting and relative orientations}
A constant pair $C_x$ joins the gauge critical point $a_x$
to the symplectic intersection point $x$. Its operator is
invertible. Let $L_S(x)=L_G(a_x)$. For a mixed configuration
$c$ with gauge input $a$ and symplectic output $x$, shifting
and gluing identify it, through a path of configurations,
with a configuration of the form $A\# C_x$. Here $A$ is a
gauge path from $a$ to an appropriate lift of $a_x$.

\begin{proposition}\label{prop:transport}
Under Hypothesis~\ref{hyp:regularity}, there is a relative
orientation isomorphism
\begin{equation}\label{eq:Gamma}
 \Gamma_c:L_G(a)\otimes\det D_c\longrightarrow L_S(x),
\end{equation}
defined up to positive scalar, which agrees with gauge
orientation gluing at $A\#C_x$. It is independent of the
shifting path, natural under the lifting group, and coherent
under ordinary gauge and strip gluing.
\end{proposition}
\begin{proof}
First consider a loop of configurations with fixed limits.
The relative family shifting in
Hypothesis~\ref{hyp:regularity}(c) deforms it to a loop of
symplectically constant configurations. A common gluing
length works for the compact family. At that end of the
deformation, the determinant bundle is the gauge
orientation bundle tensored with the determinant of the
invertible constant pair. The gauge bundle is oriented by
the almost-complex construction, and the second factor has
its canonical orientation. The orientation monodromy of
the loop is therefore trivial. This proves path independence
of the orientation transported from the constant reference.

For a strip attached at the output, first shift the original
mixed configuration to $A\#C_x$. Then shift $C_x$ together
with the attached strip to another gauge path followed by
$C_y$. Associativity of determinant gluing identifies both
orders of this construction with gauge gluing at the final
reference. Continuity transports that equality back to the
original configuration. The gauge-end statement is the same
argument with the extra path attached before shifting.

The lifting automorphism preserves gauge orientations,
conjugates the constant-pair operator, and commutes with
gluing. Conjugation preserves the canonical orientation of
an invertible operator. Hence the construction descends to
the framed quotient. If shifting reaches the other lift of
$a_x$, the two end lines are identified by this same action.
\end{proof}

The pure strip orientation is defined by applying
\eqref{eq:Gamma} to $C_x\#u$ and cancelling the constant
pair. This is the transported orientation construction
of~\cite[Section~4.7]{DFL}. For a mixed triangle there are
several gauge inputs; their lines are listed in cyclic
order and the same construction is performed on all arms.
Associativity is applied before any factor is reordered.

\subsection{Parameters and boundary orientations}
A parameterized family is oriented by placing the oriented
parameter space before the determinant line of the vertical
operator. At a regular point this induces the usual
orientation of the finite-dimensional zero set. To see
continuity when the vertical kernel dimension changes,
choose one stabilization on a neighborhood in the parameter
space. The orientation is then the exterior-power
orientation of the stabilized kernel with the fixed
stabilization factor cancelled. Transition maps between
such neighborhoods are the determinant-bundle transition
maps; no choice of a basis of the moving kernel is required.

For a one-parameter regular family, the boundary orientation
is outward normal first. At a strip break the translation
vector is the positive vector $\partial_su$. Moving its
odd line past the preceding generator factors gives
exactly the signs in~\eqref{eq:tensor}. At a parameter
endpoint the two outward normals are opposite. It follows
that a compactified one-dimensional family yields an
identity of the form
\begin{equation}\label{eq:family-boundary}
 F_1-F_0=\delta K,
\end{equation}
with the Koszul signs already included in $\delta$.
If an orientation for the whole family has been chosen
only up to sign, the corresponding endpoint identity
has one overall sign. It never has independently chosen
signs for its separate matrix entries.

The compatibility of a geometric boundary collar with
this determinant-line orientation is part of the gluing
assertions in the analytic hypotheses. Merely orienting
the interior of a moduli space would not prove
\eqref{eq:family-boundary}.

\section{The mixed isomorphism over the integers}\label{sec:vertex}

\subsection{Differentials and the vertex map}
For an index-one gauge trajectory or holomorphic strip,
its sign is positive when the relative orientation agrees
with its translation vector. For a regular index-zero
mixed solution, its sign is the comparison of
\eqref{eq:Gamma} with the canonical element of its
invertible determinant line. Set
\begin{align}
 d_Go_a&=\sum_b\sum_{[A]\in\mathcal M_G^0(a,b)}
                   \operatorname{sign}(A)o_b,\notag\\
 d_So_x&=\sum_y\sum_{[u]\in\mathcal M_S^0(x,y)}
                   \operatorname{sign}(u)o_y,\label{eq:differentials}\\
 N_po_a&=\sum_x\sum_{[c]\in\mathcal M_N^0(a,x)}
                   \operatorname{sign}(c)o_x.\notag
\end{align}
The superscript $0$ denotes dimension after quotienting
translation when that action is present. All sums are
finite by the compactness and energy-index bounds of
Paper~I.

\begin{proposition}\label{prop:N-chain}
The transported orientations give
$d_S^2=0$ and $d_SN_p=N_pd_G$ over $\ZZ$.
\end{proposition}
\begin{proof}
The compactified one-dimensional strip space has two-level
broken strips as its boundary. Coherence in
Proposition~\ref{prop:transport} identifies their boundary
orientation with the product of the two zero-dimensional
signs, with the sign from the translation factor. The
oriented boundary count is zero, yielding $d_S^2=0$.
For the mixed vertex, the two types of boundary are a
gauge trajectory followed by a mixed solution and a mixed
solution followed by a strip. The outward-normal
orientations of the two ends give opposite signs. Their
sum is $d_SN_p-N_pd_G$. The index and energy restrictions
exclude additional boundary strata in this dimension.
\end{proof}

\subsection{The unit diagonal}
The filtration used in Paper~I orders the finite generator
set so that every index-zero mixed solution goes from a
generator to a generator of no larger filtration. At equal
filtration the only solutions are the unique small
deformations of the constant pairs. This is a statement
about the set of solutions, before orienting it.

\begin{lemma}\label{lem:triangular}
Let $F$ be an integral endomorphism of a finite filtered
free module. Suppose its associated graded map is diagonal
with each diagonal entry in $\{1,-1\}$. Then $F$ is an
integral isomorphism. Its inverse preserves the filtration.
\end{lemma}
\begin{proof}
Write $F=D(1+U)$, with $D$ the diagonal sign matrix and
$U$ strictly filtration-lowering. If there are $r$
filtration steps then $U^r=0$, and
\[
 F^{-1}=(1-U+U^2-\cdots+(-1)^{r-1}U^{r-1})D^{-1}.
\]
All coefficients are integral.
\end{proof}

Each diagonal solution contributes $1$ or $-1$, since it
is one regular point. Thus the invertibility of $N_p$
already follows from Lemma~\ref{lem:triangular}; it does
not require first proving that every diagonal entry is
positive. With the normalization by constant pairs, the
entries are in fact $+1$.

\begin{proposition}\label{prop:N-isomorphism}
The map $N_p$ is a chain isomorphism. Its diagonal entries
are $+1$ in the constant-pair normalization. Its homotopy
class is independent of regularizing data, after the
corresponding gauge continuation is inserted.
\end{proposition}
\begin{proof}
Take the segment from zero regularizing data to the
chosen small data. At a constant pair the linearization
is invertible. The chart implicit-function theorem
therefore supplies a nearby branch of regular solutions.
The uniqueness statement in the triangularity argument
identifies this branch with the diagonal solution. The
fixed-center regularity in
Hypothesis~\ref{hyp:regularity} makes the branch continuous
for determinant-line transport. Its sign is locally
constant and begins at the canonical sign $+1$.

The filtration now gives an inverse integral matrix.
Since $N_p$ commutes with the differentials, its inverse
does too. Finally connect two choices of regularizing
data by a generic small path, allowing the gauge end
perturbation to vary and retaining the fixed primary
data. The compactified one-dimensional mixed spaces give
\eqref{eq:family-boundary}. The gauge continuation accounts
for the changing gauge complex, while the strip end is
unchanged. This proves the stated independence.
\end{proof}

The secondary perturbations are supported away from the
strip end, so their variation does not change the strip
operator. The segment is chosen in a star-shaped
smallness neighborhood where all the reference
constant-pair operators remain invertible. This is why
an arbitrary path of perturbations cannot be substituted
in the normalization argument.

\subsection{Parametric orientations}
The same reasoning applies to compact parameter families
of primary and geometric data whenever their asymptotic
intersections remain transverse. Locally choose
fixed-center charts and stabilizations. A nearby regular
solution has small coordinates in such a chart. The
line orientation pulled back from
Proposition~\ref{prop:transport} varies continuously,
while the canonical orientation of an invertible
operator is continuous on the invertible locus. Their
ratio is consequently locally constant. This fact will
be used in the collapse and closed-cut comparisons.
It does not require identifying the topology of every
Sobolev quotient with the smooth quotient globally.

\section{Triangles and orientations under collapse}\label{sec:collapse}

\subsection{The triangle product}
Order three Lagrangians $L_0,L_1,L_2$ by the boundary
orientation of the triangle. Write $C_{ij}=\CF(L_i,L_j)$.
Transport orientations through the mixed triangle and
put all generator lines before its determinant line.
The resulting product
\[
 \mu:C_{01}\otimes_R C_{12}\longrightarrow C_{02}
\]
satisfies
\begin{equation}\label{eq:leibniz}
 d\mu(a,b)=(-1)^{|b|}\mu(da,b)+\mu(a,db).
\end{equation}
Indeed the first strip break requires its translation
line to pass the second input generator; the second
break does not. This is the boundary calculation for
the one-dimensional triangle space.

If $\theta\in C_{12}$ is a homogeneous cycle, then
$R_\theta(a)=\mu(a,\theta)$ is a closed map of degree
$|\theta|$. The product itself is even in this convention.
For a change of gauge base at any of the three vertices,
the corresponding symplectic conversion is diagonal on
the same generator lines. Applying the conversions with
\eqref{eq:tensor-maps} gives the product in the new bases;
no unsigned tensor of two possibly odd conversions is
used.

\subsection{The determinant comparison at the collapsed limit}
Let $V$ be a regular index-zero limiting mixed solution
in the collapse theorem of Paper~I. Denote its approximate
solution on the domain of collapse scale $\epsilon$ by
$W_\epsilon(V)$ and the actual solution by
$\mathcal G_\epsilon(V)$. The augmented operator estimates
in Paper~I give the inverse on the orthogonal complement
of the stationary coordinates and retain the stationary
finite-dimensional factor explicitly. For the regular
index-zero solution the complete realized operator is
invertible.

The determinant comparison is constructed using the
same finite-dimensional stabilization on the limiting
and preglued problems. The patched kernel map identifies
the two stabilized kernels for small $\epsilon$.
Cancelling the stabilization factors gives
\begin{equation}\label{eq:collapse-det}
 \det D_V\longrightarrow\det L_{W_\epsilon(V)}.
\end{equation}
At a constant reference configuration this is the
positive identity. Coherence under gauge and strip
gluing propagates the normalization to each component
of the mixed configuration space, by the shifting
argument of Proposition~\ref{prop:transport}.

It remains to compare the preglued operator with the
operator at the actual solution. A potential difficulty
is that the chart used for the contraction estimate
need not be the chart used to define the determinant
bundle. The second derivative of an arbitrary fixed
mixed chart need not have a uniform bound in the
collapse norm. We use the following argument instead.

\begin{lemma}[Transport around a square]\label{lem:square}
Let $D_{s,t}$ be a continuous family of Fredholm operators
on a square. Suppose the operators on three edges are
invertible. Parallel transport of orientations on the
fourth edge carries the canonical orientation at its
first endpoint to the canonical orientation at its
second endpoint.
\end{lemma}
\begin{proof}
The orientation double cover of the determinant bundle
is trivial over the square. Transport around its
boundary is therefore the identity. On each of the
three invertible edges the canonical orientation is a
continuous section, so transport is positive. The
remaining edge must have the same property at its
endpoints.
\end{proof}

\begin{proposition}\label{prop:collapse-sign}
The map~\eqref{eq:collapse-det}, followed by transport
from $W_\epsilon(V)$ to $\mathcal G_\epsilon(V)$, preserves
the sign of the zero-dimensional solution.
\end{proposition}
\begin{proof}
Use the contraction chart of Paper~I on one horizontal
edge of a square. The straight coordinate segment from
the preglue to the solution stays in the region where
the derivative differs from the invertible approximate
operator by a sufficiently small bounded operator.
Consequently this edge is invertible. At the solution,
change from the contraction realization to the
admissible orientation realization by
Lemma~\ref{lem:realization}; this vertical edge is
invertible. At the approximate configuration use the
same source and target identifications as those defining
\eqref{eq:collapse-det}; the comparison gives the other
invertible vertical edge.

The family in the interior of the square consists of
admissible Fredholm realizations. This uses
Hypothesis~\ref{hyp:regularity} and the exact matching
chart, including the conformal gauge row. The fourth
edge is the path of orientation operators at the
configurations between the preglue and the solution.
We need not assert that this entire edge is invertible.
Lemma~\ref{lem:square} nevertheless transports the
canonical orientation between its endpoints. Combining
this with~\eqref{eq:collapse-det} proves the claim.
\end{proof}

If the conformal metric in the collapse estimate is
changed within the unperturbed collar class, the zero
set of the mixed equation is unchanged. At a solution
Lemma~\ref{lem:realization} compares the gauge rows by
invertible paths, so the sign is unchanged. This permits
use of the metric of Paper~I without introducing a new
orientation convention.

The construction is compatible with compact parameter
families because its stabilizations and operator
bounds may be chosen locally uniformly. Associativity
of ordinary gluing identifies the determinant maps at
broken ends. This proves Theorem~\ref{thm:B}.

\subsection{Vertex identifications and two-handles}
In the triangle degeneration a completing vertex is
identified with a product of a gauge path and a
constant mixed pair. Normalize this identification
using the corresponding end frames, with the outgoing
coordinate translated in the positive direction on
every arm. The constant pair has sign $+1$ and the
gauge path uses the fixed object line. The resulting
vertex identification has sign $+1$ at the reference
configuration. Shifting and gluing show that it has
that sign throughout its configuration component.

For the standard surgery triple let $\theta_G$ denote
the distinguished gauge cycle in the auxiliary
connected sum of $S^1\times S^2$'s. Normalize its
symplectic counterpart by
\[
 \theta_S=N_{12}\theta_G.
\]
The unit diagonal and the usual filtration of that
auxiliary complex identify this with the distinguished
top cycle, with the chosen sign. The mixed triangle
family has one endpoint equal to a gauge triangle
followed by the mixed map, and the other endpoint
is a symplectic triangle with the two mixed inputs.
Proposition~\ref{prop:collapse-sign} compares their
orientations. The broken boundary terms give a
homotopy, hence
\begin{equation}\label{eq:triangle-square}
 N_{02}\mu_G\simeq
 \pm\mu_S(N_{01}\otimes_R N_{12}).
\end{equation}
Evaluating the second input on $\theta_G$ gives
\begin{equation}\label{eq:twohandle}
 N_{p_1}\IG(W_2)\simeq\pm S_{W_2}N_{p_0}.
\end{equation}
Changing bases conjugates this identity by the
right-placed conversions and uses
\eqref{eq:compose-homotopy}. Twisting the bundle on
the two-handle changes the gauge triangle class but
not the orientation argument; its relative bundle
is retained throughout the mixed family.

For simultaneous gluing at two triangle ends, use
one stabilization off both ends. The determinant
comparison is independent of the order of gluing,
up to the graded reordering already present in
\eqref{eq:tensor-maps}. Thus the triangle systems
are coherent at corners, which is needed when the
same product is later folded into a compression cup.

\section{Torus excision over the integers}\label{sec:excision}

\subsection{The closed cut and its orientation data}
The folded triangle has a closed gauge cut
\[
 Z\cong(-M)\#M^+\#T^3.
\]
The odd tori $T_a,T_b\subset Z$ are homologous, and the
lifting group at this cut is generated by
$s_Z=\PD[T_a]\pmod2$. Excision along these tori gives
\[
 Z'=(-M)^\#\sqcup(M^+)^\#.
\]
Write $\mathsf E:C_G(Z)\to C_G(Z')$ for the gauge
excision map and $\overline{\mathsf E}$ for its reverse.
The tensor complex at $Z'$ uses the right convention.

The almost-complex structure induced by the folded
triangle at $Z$ need not equal the one induced by the
mixed torus domain. A conversion at this interior cut
is therefore required. We choose the structures so that
\begin{equation}\label{eq:cut-Chern}
 c_1(J)[T_a]\equiv0\pmod4.
\end{equation}
For two such structures, the character
\eqref{eq:gauge-character} on $s_Z$ is trivial.
Consequently the conversion at $Z$ is defined on the
quotient complex.

Here is why the choice is available. On an oriented
four-manifold with boundary, spin$^c$ structures form
an affine space for $H^2$, and their first Chern classes
change by twice that class. If the image of
$H^2(B;\ZZ)\to H^2(T_a;\ZZ)$ is $m\ZZ$, write
$[T_a]=m\sigma$ modulo torsion. When $m$ is odd, twisting
changes half the Chern evaluation modulo two and permits
\eqref{eq:cut-Chern}. When $m$ is even, the evaluation
on $\sigma$ is even: the restriction of $w_2(B)$ to
its oriented three-dimensional boundary vanishes, and
$\sigma$ may be taken there. Thus the evaluation on
$T_a$ is already divisible by four. The almost-complex
structure is then made cylindrical near $Z$. This
argument is applied with the fixed torus germ away
from the interior cut.

\begin{convention}\label{conv:interior-line}
The determinant line of the conversion at an interior
cut is placed after the complete input generator line
and before the complete configuration determinant line.
Its position is fixed relative to the entire
configuration, independently of where the geometric
cut is drawn.
\end{convention}

Suppose the configuration is written as a sequence of
glued pieces. Moving an odd conversion line with the
cut through that sequence would move it past the
determinant line of the intervening piece. The sign
would then depend on its index. Convention
\ref{conv:interior-line} avoids this dependence: moving
the line to a fixed reference position costs a sign
determined by the entire configuration, which is
unchanged by moving the cut. At index zero the
alternative fixed placement after the complete
configuration line agrees with this one.

\subsection{Lifting groups and gluing multiplicities}
It is useful to perform the gauge excision calculation
before passing to a quotient. Let $\mathcal G_0(X)$
be the subgroup of bundle automorphisms with trivial
lifting class. If a finite group $\phi_X$ acts freely
and orientation preservingly on the relevant
zero-dimensional solution set, its quotient count is
\begin{equation}\label{eq:orbit-count}
 I(X)^{\phi_X}(\bar a_0,\bar a_1)
 =\frac{1}{|\phi_X|}
   \sum_{a_i\in\bar a_i}\#\mathcal M(X;a_0,a_1).
\end{equation}
The numerator is divisible by $|\phi_X|$: it is a
sum of free orbits, each with a constant sign.

\begin{lemma}\label{lem:gluing-multiplicity}
Let $X$ be connected and let $Y$ be a closed
hypersurface with cut manifold $X_Y$. Suppose the
critical limits on $Y$ have trivial stabilizer in the
adjoint determinant-one gauge group. For each matched
configuration on $X_Y$ that glues, the resulting
$\mathcal G_0(X)$-classes form a torsor for
\[
 K_Y=\ker\bigl(H^1(X;\ZZ/2)\longrightarrow
                         H^1(X_Y;\ZZ/2)\bigr).
\]
In particular their multiplicity is $|K_Y|$.
\end{lemma}
\begin{proof}
Choose representatives whose two limits agree under
the fixed gluing identification. If two choices glue
to the same determinant-one class on the cut manifold,
their difference is an automorphism whose lifting
class restricts to zero on $X_Y$. Conversely every
such class changes the gluing while leaving the cut
class fixed. The stabilizer condition makes this action
free. Its transitivity follows by comparing the two
gluing identifications on the two copies of $Y$.
The Mayer--Vietoris sequence identifies $K_Y$ with the
image of the connecting map from $H^0(Y;\ZZ/2)$.
The analytic assertion that every required matched
solution glues and that there are no others is
Hypothesis~\ref{hyp:selfgluing}.
\end{proof}

A finite lifting group acts freely on the solutions
in the present calculation. Indeed, a fixed solution
would give a parallel automorphism. Its restriction
to an irreducible end must be the identity once its
lifting class there vanishes. A parallel automorphism
which is the identity on an end is the identity
everywhere on a connected cobordism. The almost-complex
orientation construction preserves signs under the
whole group. Thus~\eqref{eq:orbit-count} applies.

\subsection{The capped comparison}
Let $V=\mathcal E\cup_{Z'}\overline{\mathcal E}$ be
the composite of excision and reverse excision.
The standard local model is $T^2$ times the doubled
excision polygon. It contains a nonseparating
hypersurface $\Sigma=T^2\times\gamma$. Cut along
$\Sigma$ to obtain $V^\circ$, and cap the two new
boundary components by $T^2\times D^2$. The capped
cobordism $V'$ is a product on $Z$; this is the
geometric comparison in torus excision
\cite[Theorem~5.6]{KM-unknot}.

The determinant-one complex on $\Sigma$ has two
critical points $g_1,g_2$. The class
$a=\PD_\Sigma[T^2]$ interchanges them. They must be
kept separate when the caps are attached: the caps
carry the trivial lifting group, since $a$ does not
extend across their meridional disks.

Write $s=\PD_Z[T_a]\pmod2$. In $H^1(V^\circ;\ZZ/2)$
let $p$ be its pullback and $h_1,h_2$ the meridional
classes of the two removed cap neighborhoods. The
restrictions are
\[
\begin{array}{c|cccc}
 &Z_{\rm in}&Z_{\rm out}&\Sigma_{\rm in}&\Sigma_{\rm out}\\\hline
 p&s&s&0&0\\
 h_1&s&0&a&0\\
 h_2&s&0&0&a\\
 h_1+h_2&0&0&a&a
\end{array}
\]
The group $H^\circ=\langle p,h_1,h_2\rangle$ has order
$8$. It contains both groups needed on the cut
cobordism, but is not itself the pullback of either
of the groups on $V$ and $V'$.

More precisely, $\phi_V$ has order $8$, its image on
$V^\circ$ is $\langle p,h_1+h_2\rangle$ of order $4$,
and its kernel is $K_\Sigma$ of order $2$. The capped
product has group $\phi_{V'}=\langle p\rangle$ of
order $2$. To check these assertions, the two sides
of the excision polygon give classes $P_+,P_-$ and
the disconnected intermediate cut gives a class $c$.
Their restrictions are $p,p,h_1+h_2$, respectively.
They are independent on $V$, as seen by loops
crossing each of the two polygon sides and the
intermediate cut. Their sum $P_++P_-$ is dual to
$\Sigma$, generates the restriction kernel, and is
nonzero because $\Sigma$ is nonseparating.

Choose invariant regularizing data on $V^\circ$.
Let
\[
 D(\alpha,\sigma;\beta,\tau)
 =\#\mathcal M(V^\circ;(\alpha,\sigma),(\beta,\tau)),
\]
where $\sigma,\tau\in\{g_1,g_2\}$. Put
$\Delta(\alpha,\beta)=D(\alpha,g_1;\beta,g_1)$.
The action of $h_1+h_2$ makes the two diagonal
blocks equal; the action of $h_2$ identifies the
off-diagonal block with $\Delta(s\alpha,\beta)$.
Thus, after summing over endpoint lifts, diagonal
and off-diagonal blocks have the same sum.

Each cap has one regular flat solution in its
specified bundle and therefore invariant $\pm g_1$
(or the corresponding dual generator). The odd
bundle on $T^2$ has no infinitesimal flat deformations,
which gives regularity. The index-energy formula
excludes a nonflat solution of index zero. Hence
the cap coefficient is a unit over $\ZZ$.

\begin{proposition}\label{prop:capped}
Under Hypothesis~\ref{hyp:selfgluing}, the quotient
counts for $V$ and $V'$ have the same absolute
multiplicity. In each fixed degree they agree up
to the determinant-gluing sign.
\end{proposition}
\begin{proof}
For $V$, self-gluing has multiplicity $2$ and there
are two equal diagonal blocks. Quotienting by the
group of order $8$ gives
\[
 \frac{2\cdot2}{8}
   \sum_{\alpha\in\bar a,\,\beta\in\bar b}
       \Delta(\alpha,\beta).
\]
For $V'$, the caps have unit coefficients, their
separating gluings have multiplicity $1$, and the
group has order $2$. The result is
\[
 \frac12
   \sum_{\alpha\in\bar a,\,\beta\in\bar b}
       \Delta(\alpha,\beta).
\]
If the second cap selects $g_2$, the symmetry of
the off-diagonal block gives the same sum. The
orientation signs are those in the gluing
hypothesis. All group actions preserve the
orientations, so no cancellation occurs inside
an orbit. The two numerical factors agree.
\end{proof}

In particular no division by two is performed on
a Floer homology group. The fractions above express
integer orbit counts. Computing the caps only
in the rank-one quotient would lose the distinction
between matching two lifts and matching two orbits.

The one-dimensional gluing family identifies these
counts with the composite excision map. The same
argument for the other composite gives the reverse
comparison. Consequently $H(\mathsf E)$ and
$H(\overline{\mathsf E})$ are isomorphisms. On a
homogeneous class their composite is that class
up to a sign depending on its degree. This is the
only inverse-excision statement needed for the
handle state below. By Lemma~\ref{lem:acyclic}, the
closed excision maps are homotopy equivalences
(after shifting degrees if necessary). We do not
replace a possibly degree-dependent sign by a
single scalar on an ungraded complex.

\subsection{The mixed excision identity}
Let $N_Z$ be the mixed map on the closed-cut
configuration, and $N_-,N_+$ the vertex maps on
the two separated ends. The symplectic output of
$N_Z$ is a Floer complex in a product target.
Folding identifies its generators with pairs
$(x,y)$ in the reflected source and target.
The identification of their orientation lines is
denoted by $\mathcal K$; it is fixed by comparison
at split constant reference configurations. If
an input conversion has parity $t_{\rm in}$,
write $\mathcal K_R=\mathcal K\eps^{t_{\rm in}}$
for its closed-map convention.

\begin{theorem}[Mixed excision, conditional]\label{thm:mixed-excision}
Under the analytic hypotheses, with conversions at
the closed cut placed as in Convention
\ref{conv:interior-line}, one has
\begin{equation}\label{eq:excision-reverse}
 \mathcal K_RN_Z\overline{\mathsf E}
       \simeq\pm(N_-\otimes_R N_+).
\end{equation}
For a homogeneous cycle $z\in C_G(Z)$ this implies
\begin{equation}\label{eq:excision-cycle}
 [\mathcal K_RN_Zz]
 =\pm[(N_-\otimes_R N_+)\mathsf E z].
\end{equation}
All terms are expressed at the fixed end bases.
\end{theorem}
\begin{proof}
Use the closed-cut gluing from Paper~I to express
$N_Z\overline{\mathsf E}$ as the joined mixed
count. The ordinary neck homotopy gives a signed
chain homotopy by Lemma~\ref{lem:line-gluing} and
\eqref{eq:family-boundary}. Next vary the width
and complete the torus ends. This is exactly the
comparison asserted in
Hypothesis~\ref{hyp:completion}. At the completed
limit the zero-dimensional solutions are pairs
of solutions on the two separate domains.

At split reference configurations the determinant
identification is the product identification by
Hypothesis~\ref{hyp:completion}(b). Coherence at
gauge and strip ends and the shifting argument
propagate this comparison to every component.
Thus the only generator-dependent factor is the
fixed folding identification $\mathcal K$ itself;
there is no additional uncontrolled diagonal
sign. Listing the two output lines in order gives
$N_-\otimes_R N_+$. The comparison families supply
the homotopy in~\eqref{eq:excision-reverse}.

Apply~\eqref{eq:excision-reverse} to
$\mathsf E z$. The capped comparison says
$[\overline{\mathsf E}\mathsf E z]=\pm[z]$
for this homogeneous class. Applying $N_Z$ gives
\eqref{eq:excision-cycle}.
\end{proof}

Equality~\eqref{eq:excision-cycle} is equality
of homology classes of tensors. If the tensors
differ by a boundary, applying the chain
isomorphism~\eqref{eq:beta} produces a boundary
in the Hom complex. This is the mechanism that
will turn the handle-state calculation into a
chain homotopy. No inference from equality of
induced maps on homology is involved.

\section{The compression cup before reflection}\label{sec:cup}

\subsection{The tensor supplied by the folded triangle}
Let $C$ be the compression cobordism for $k$
one-handles, let $p$ be the source presentation,
and put $q=C\circ p$. The folded target is a
product of a negatively oriented source moduli
space and the target moduli space. The graph of
restriction over $C$ is a Lagrangian
correspondence in this product. Two auxiliary
height Lagrangians supply distinguished
intersection cycles $q_0^*,q_1^*$.

The folded triangle product, with the conversion
at $Z$ inserted, gives a cycle
\begin{equation}\label{eq:cup-tensor}
 z_C=\mathcal K_R\mu_X^\Theta(q_0^*\otimes_R q_1^*)
       \in\CF(\bar p)\otimes_R\CF(q).
\end{equation}
The triangle and the folding are defined by
holomorphic curves and their transported
orientation lines. Define
\begin{equation}\label{eq:check-cup}
 \check U_C=B_{\CF(\bar p)}(z_C):
       \CF(\bar p)^\vee\longrightarrow\CF(q).
\end{equation}
For $k=0$ denote this map by $B(p,r)$, where
$r$ is the upper height presentation. These
maps have not yet used a symplectic reflection
pairing.

The top cycles are normalized using the integral
mixed maps on the auxiliary vertices:
$q_i^*=\pm N_i\theta_i$. The sign is one sign
of each distinguished cycle. The auxiliary
complexes have the standard filtration, and the
unit-diagonal calculation in
Proposition~\ref{prop:N-isomorphism} makes this
an integral normalization, not merely a
mod-two one.

\subsection{The gauge handle state}
Let $B_k$ be the gauge filling whose closed
boundary is $Z$, obtained from the folded
triangle and its two auxiliary fillings. Its
relative invariant is a homogeneous cycle
$\Theta_{B_k}\in C_G(Z)$. Choose a single
almost-complex structure on this filling and
use its restrictions on the pieces. Convert
at $Z$ to the structure used for mixed torus
excision. After excision the resulting state
is a cycle on
$C_G(-M)\otimes_R C_G(M^+)$.

Bending the source end of the one-handle
cobordism $W_k$ gives the identity
\begin{equation}\label{eq:handle-state}
 (\Phi_{G,M}\otimes_R1)[\mathsf E\Theta_{B_k}]
       =\pm[\ten(\IG(W_k))].
\end{equation}
To check the sign convention, start with the
product cylinder: its bent state is the
coevaluation tensor. Glue $W_k$ to its second
end and use the ordered determinant gluing.
This gives the tensor of the right-placed
map. An odd endpoint conversion is already
part of that map through
\eqref{eq:theta-grading}; moving it to the
other side would change the tensor by a
grading operator. The residual signs in
\eqref{eq:handle-state} are therefore scalar.
The same argument holds for $k=0$, when the
cobordism map is gauge continuation.

\subsection{The cup formula}
Define the transposed reflected mixed map
\begin{equation}\label{eq:n-p}
 n_p=\ev_{C_G(M)}^{-1}
       (N_{\bar p}\Phi_{G,M}^{-1})^\vee:
       \CF(\bar p)^\vee\longrightarrow C_G(M).
\end{equation}
It is an even chain isomorphism, using the
paired choices of gauge base. This definition
uses gauge duality and the mixed map at the
reflected presentation; it makes no assumption
about symplectic duality.

\begin{proposition}[Cup formula]\label{prop:cup-formula}
Under the analytic hypotheses,
\begin{equation}\label{eq:cup-formula}
 \check U_C\simeq\pm N_q\IG(W_k)n_p.
\end{equation}
In particular, for a lower datum $p$ and a
compatible upper height datum $r$,
\begin{equation}\label{eq:cup-zero}
 B(p,r)\simeq\pm N_r\cont_G(p,r)n_p.
\end{equation}
\end{proposition}
\begin{proof}
Apply the integral triangle comparison to
the two distinguished input cycles. It
identifies the folded triangle output,
up to a boundary and one sign, with
$N_Z\Theta_{B_k}$, using the conversion
at $Z$ on both sides. Apply
\eqref{eq:excision-cycle} to the homogeneous
cycle $\Theta_{B_k}$. The result is
\[
 [z_C]=\pm[(N_{\bar p}\otimes_R N_q)
                          \mathsf E\Theta_{B_k}].
\]
Equation~\eqref{eq:handle-state} identifies
this with the tensor of $N_q\IG(W_k)n_p$.
The identification follows directly from
\eqref{eq:beta} and the graded double-dual
map in~\eqref{eq:n-p}. Since the difference
of the tensors is a boundary, their images
under $B_{\CF(\bar p)}$ differ by $\delta K$.
This proves~\eqref{eq:cup-formula}. The same
calculation with $k=0$ gives
\eqref{eq:cup-zero}.
\end{proof}

This proves a formula for the cup tensor
without selecting the reflection pairing
needed for the source of Paper~I's map.
The next sections construct that pairing
and show that it gives the ordinary
symplectic continuation when $k=0$.

\section{Sphere characters and variation of primary data}
\label{sec:characters}

\subsection{The character of an orientation system}
For a transverse pair $P=(P_0,P_1)$ let
$\mathcal P(P_0,P_1)$ be the path space of paths in
$\Msf_h$ starting on $P_0$ and ending on $P_1$.
For $h\ge2$, the compression-body Lagrangians are
simply connected and have vanishing second homotopy
group. The relative homotopy classes of strip
configurations with fixed ends are consequently a
torsor for $\pi_2(\Msf_h)\cong\ZZ$ in the lift
convention of~\cite[Section~2]{DFL}. Fix its positive
generator $S_h$; attaching it changes the index by
four. For $h=1$ the moduli space is a point and the
sphere-character assertions below are empty.

The determinant line of a sphere operator, together
with the inverse determinant of the matching tangent
space at its node, has its complex orientation.
Compare node gluing with the coherent strip
orientation $\Gamma^P$. The ratio is a sign
\begin{equation}\label{eq:sphere-character}
 \chi^P:\pi_2(\Msf_h)\longrightarrow\{1,-1\}.
\end{equation}
Associativity of node gluing makes it a character.
Multiplying generator lines by signs does not change
it: the two endpoint factors occur in both sides of
the comparison. Multiplying them by their grading
signs has the same property. Thus the character is
information about the configuration-level coherent
system, beyond the signs of its differential.

At a transverse point $x$, a second description is
available. Attach a fixed representative sphere to
the constant mixed pair $C_x$ and compare its
transported mixed orientation with the product of
the canonical orientation of $C_x$ and the complex
node orientation. Denote the ratio by
$\chi(h,x;d)$, where $h$ is the primary datum and
$d$ contains the regularizing and almost-complex
data.

\begin{lemma}\label{lem:pointed-character}
At regular data,
$\chi(h,x;d)=\chi^P(S_h)$ for every transverse
intersection point $x$. The ratio is constant along
any compact smooth family of primary data carrying
a persistent transverse intersection point. It is
also independent of small regularizing data at
fixed primary data and fixed gauge base.
\end{lemma}
\begin{proof}
Glue a strip to the constant pair and attach the
sphere away from the ordinary gluing collar. The
two possible orders of these operations induce
the same determinant-line isomorphism. The sphere
and node factors are complex, hence even as real
graded lines, so commuting them with the strip
factor gives no sign. Cancelling the strip
orientation gives equality of the pointed and
strip definitions. A strip configuration exists
between any two intersection points, in some
relative homotopy class; it need not solve the
holomorphic equation. Applying the same argument
to that configuration shows independence of $x$.

For a family with a persistent point, the
constant-pair operators, the transported
orientation, and the node-gluing map vary
continuously by Hypothesis~\ref{hyp:regularity}.
The ratio is a continuous function to
$\{1,-1\}$ and is therefore constant. To vary
regularizing data, use the segment to zero data.
The perturbations are supported off the strip
end, and the constant-pair operators stay
invertible. Transporting their canonical
orientations along the segment proves the last
assertion. Compatible almost-complex structures
can be varied through the same argument, with
the gauge base held fixed.
\end{proof}

The persistent point need only remain transverse
at that point. Other intersections of the two
Lagrangians may become degenerate. The lemma is
about determinant lines over configurations and
therefore does not require a regular moduli
space along the family.

\subsection{A common ball of primary perturbations}
Let $\mathcal C_i$ be the finite real span of the
sign-sensitive holonomy functions on $C_i^\circ$,
modulo constants. Use the norm
\begin{align}
 \|h_i\|_i={}&\sup_B\|\nabla_Bh_i\|_{L^\infty}
  +\sup_{B\ne B'}
   \frac{\|\nabla_Bh_i-\nabla_{B'}h_i\|_{L^2}}
        {\|B-B'\|_{L^2}},\label{eq:primary-norm}\\
 \|h\|={}&\max(\|h_0\|_0,\|h_1\|_1).\notag
\end{align}
The suprema are taken over smooth connections in a
fixed determinant bundle. Cylinder functions have
finite norm by the holonomy estimates of
\cite[Proposition~6.1]{DFL}. If the norm vanishes,
the gradient vanishes on the affine connection
space, so the function is constant; thus this is
a norm modulo constants.

Choose a finite trace family $F_0$ whose restrictions
to the compact flat-connection moduli spaces give
embedding coordinates. In the model
$L(C_i)\cong\SU(2)^{h-1}$, quaternion coordinates
can be obtained from sign-sensitive traces using
the fixed holonomies on the odd torus. Such traces
distinguish central sign changes; passing to
unsigned $\SO(3)$ traces would lose that property.
For a finite family $F$ containing $F_0$, put
\[
 \mathcal U(F)=\{h\in\operatorname{span}F:\|h\|<r_*\},
 \qquad \mathfrak B=\bigcup_F\mathcal U(F).
\]
The radius is independent of the additional
functions in $F$, for the fixed splitting and
metric data. These sets are nested under
inclusion of finite families and are convex.

\begin{proposition}\label{prop:small-ball}
The radius $r_*>0$ can be chosen so that every
$h$ in this ball has the following properties.
\begin{enumerate}[label=(\roman*)]
\item The perturbed flat moduli space of each
compression body is a smooth compact manifold
identified with its unperturbed moduli space;
its boundary restriction is an embedded
Lagrangian close to the unperturbed one.
\item The connection representatives and their
first derivatives in the moduli parameters
are close in $L^2_1$ and $L^2$, respectively,
uniformly as $\|h\|\to0$.
\item The functions from $F_0$ remain embedding
coordinates on these perturbed moduli spaces.
\item If the two boundary Lagrangians intersect
transversely, the corresponding closed gauge
critical points are nondegenerate and are
exactly the lifts of those intersections.
\end{enumerate}
\end{proposition}
\begin{proof}
The norm gives
$\|F_B\|_{L^\infty}\le\|h_i\|_i$ for an
$h_i$-flat connection. If a sequence with norm
tending to zero escaped every gauge neighborhood
of the flat moduli space, Uhlenbeck compactness
on the compact compression body would give a
subsequence weakly converging in $W^{1,4}$,
after gauge, to a flat connection. The
$W^{1,4}$ embedding in $C^0$ is compact in
three dimensions. A Coulomb--Neumann slice
relative to the flat limit and the elliptic
estimate for $(d_B^*,d_B)$ then improve this
to convergence in $L^2_1$, a contradiction.
The boundary slice and compactness statements
are the versions for a manifold with boundary
in~\cite{Wehrheim-book}.

In a slice at a flat connection split into the
finite-dimensional harmonic space and its
complement. The linearized perturbed-flat
operator on the complement has a bounded
right inverse. The second term of
\eqref{eq:primary-norm} makes the holonomy
part a small Lipschitz perturbation in
$L^2$, so the contraction gives a graph over
the harmonic space. Differentiating its
fixed-point equation in the finite-dimensional
parameter gives the asserted first-derivative
control. Hodge decomposition and boundary
elliptic regularity identify these solutions
with the smooth perturbed-flat solutions;
the exact regularity inputs are recorded in
Appendix~\ref{app:inputs}. Compactness of the
flat moduli space makes the thresholds
uniform over finitely many slice charts.

For a trace function $\tau$, its bounded
Lipschitz gradient extends the identity
\[
 \tau(B')-\tau(B)=\int_0^1
 \langle\nabla_{B+s(B'-B)}\tau,B'-B\rangle\,ds
\]
to $L^2$ connections. Differentiating on the
graph family shows that the trace-coordinate
map remains $C^1$-close to its unperturbed
embedding. Openness of embeddings of a
compact manifold proves (iii). The matching
exact sequence for the two compression-body
boundary problems identifies the kernel
of the closed Hessian with the tangent
intersection of the Lagrangians; hence
transversality gives (iv), as in
\cite[Section~2.3]{DFL}.
\end{proof}

No radius controlling all higher derivatives
of every cylinder function is asserted here.
The norm controls the two estimates actually
used in this proposition. On a fixed
finite-dimensional perturbation family all
higher smooth seminorms are finite and are
handled separately in compact-family
regularity statements.

\subsection{Paths with only birth--death singularities}
Enlarge $F$ to include the monomials of degree at
most five in the trace coordinates $F_0$.
These remain cylinder functions: a monomial
is a smooth invariant function of the joint
finite holonomy family.

\begin{lemma}\label{lem:generic-path}
Two regular primary data in $\mathfrak B$
can be joined within a common $\mathcal U(F)$
by a smooth path with finitely many singular
parameters. At each singular parameter
there is one quadratic tangency, with
versal normal form
\begin{equation}\label{eq:birth-death}
 S_t(x,y)=x^3+\sigma(t-t_*)x+Q(y),
 \qquad\sigma\in\{1,-1\},
\end{equation}
where $Q$ is nondegenerate. All other
intersections are transverse.
\end{lemma}
\begin{proof}
Start with the segment from $h_a$ to $h_b$
and perturb it to
\[
 h_t=(1-t)h_a+th_b+\beta(t)c,
\]
where $\beta$ vanishes near the endpoints
and is positive on the interior where
singularities may occur. The segment is
compact in the open convex ball, so small
$c$ keeps the path in the ball.

Variation of the holonomy function induces
an exact Lagrangian deformation; its
Hamiltonian on the perturbed Lagrangian is
the restriction of that function to the
connection moduli space. The trace
coordinates remain embedding coordinates
by Proposition~\ref{prop:small-ball}.
Polynomials of degree at most three in
embedding coordinates realize arbitrary
third jets at one point. Polynomials of
degree at most five realize arbitrary
second jets at two distinct points. For
the latter assertion, choose an affine function $\ell$ with
$\ell(a)=0$ and $\ell(b)=1$. Given a second jet at $b$, multiply
$\ell^3$ by the quadratic Taylor polynomial at $b$ of that jet
divided by $\ell^3$. The product has degree at most five, realizes
the prescribed jet at $b$, and vanishes to order three at $a$.
Use $(1-\ell)^3$ in the same way for the jet at $a$, then add
the two polynomials. Restriction of ambient jets
to an embedded submanifold is surjective.

In a Weinstein chart, varying one
Lagrangian changes the relative generating
function by this Hamiltonian to first
order. The map on its jets is triangular
with identity diagonal. Parametric
transversality therefore excludes higher
singularities, two simultaneous tangencies,
and tangencies with zero crossing speed:
these strata have codimension at least two.
A generic one-parameter path meets only
the codimension-one birth--death stratum,
transversely. Its parameter-dependent
splitting normal form is
\eqref{eq:birth-death}.

The singular-parameter set is closed by
compactness of the Lagrangians. Each
singular parameter is isolated by the
normal form, and the endpoint neighborhoods
are regular. It is consequently finite.
No uniform separation of the two points
in the two-point jet argument is required;
transversality is applied on the open
configuration space of distinct pairs.
\end{proof}

\begin{lemma}\label{lem:persistent}
Suppose $H(C_G(M))\ne0$. The path in
Lemma~\ref{lem:generic-path} has a finite
subdivision with regular subdivision
points such that each subinterval carries
a continuous transverse intersection
point throughout the subinterval.
\end{lemma}
\begin{proof}
Every regular primary datum has at least
one intersection. Otherwise its gauge
complex would have no generators by
Proposition~\ref{prop:small-ball}(iv),
contradicting invariance and nonvanishing
of the gauge homology.

At a birth--death time the local pair
in~\eqref{eq:birth-death} is absent on
one side. If it were the only intersection
at the singular time, compactness would
force the nearby regular intersection set
on that side to be empty. Thus another
intersection exists at that time and is
transverse. The implicit function theorem
continues it across the wall.

On a closed interval with no walls, the
incidence set of intersections is a
compact one-manifold projecting locally
diffeomorphically to the interval. It is
a finite covering, hence a disjoint union
of sections. Extend the chosen transverse
point from each wall to neighboring
regular subdivision times. This gives
the required finite collection of
persistent sections.
\end{proof}

\begin{theorem}[Constancy of the sphere character]
\label{thm:character-constancy}
Fix a splitting, metric data and gauge
base, and suppose $H(C_G(M))\ne0$. Then
$\chi^P(S_h)$ is independent of the
regular datum $P$ with primary data in
$\mathfrak B$.
\end{theorem}
\begin{proof}
First remove the regularizing data at
both endpoints using
Lemma~\ref{lem:pointed-character}. Join
the primary data by the subdivided path
of Lemma~\ref{lem:persistent}. On each
subinterval the pointed character is
constant along its persistent transverse
point. At a regular subdivision time the
pointed character equals the global
character for any choice of point, so
one may switch to the persistent point
for the next subinterval. Chaining these
equalities gives constancy along the
whole path. Restore the endpoint
regularizing data. The same argument
allows variation of the compatible
symplectic almost-complex structure at
a fixed gauge base.
\end{proof}

The argument uses nonvanishing only to
find a point persisting through each
birth--death. It does not assume that
individual generators persist along an
arbitrary path, or that the regular
locus of the perturbation ball is
connected.

\section{Reflection of the transported orientations}
\label{sec:reflection}

\subsection{Reflection and sphere gluing}
The reflected presentation $\bar p$
interchanges the two compression bodies.
A strip is reflected by
\[
 \bar u(s,t)=u(-s,1-t).
\]
This map preserves the orientation of the
strip domain and keeps the surface
almost-complex structure. It reverses
the translation direction. The target
symplectic manifold retains the original
orientation; the reversed order of the
Lagrangians supplies the dual complex.
Mixed solutions themselves are not
identified by this formula. The
orientation at $\bar p$ is transported
through its own mixed domain.

For a strip from $x$ to $y$, the
reflection of its coherent orientation
has differential sign
\begin{equation}\label{eq:reflection-sign}
 \operatorname{sign}_{\rm dual}(\bar u)
   =-(-1)^{|y|}\operatorname{sign}(u).
\end{equation}
The minus sign is the reversed
translation vector; the other sign is
the graded transpose. Thus this reflected
system is the one for which the
paired generator bases give the graded
dual differential.

Attaching $S_h$ changes the index by four.
The determinant factor at a sphere node
is complex and even. The additional
transpose sign is $(-1)^{4|y|}=1$.
Consequently reflection of a coherent
strip system preserves its sphere
character. What remains is to compare
this reflected system with the system
transported through the reflected mixed
domain.

\subsection{One folded comparison}
Fix one height-fold instance with lower
primary datum $\widetilde p$ and upper
height datum $r$, both in the admissible
ball. Choose the height after the radius
of that ball, as in
Section~\ref{sec:choices}. Double the
folded triangle and attach its auxiliary
caps. The resulting disk has an incoming
strip end modeled on $\bar{\widetilde p}$
and an outgoing end modeled on $r$.
Its transported orientation is coherent
with the reflected incoming strip system
and the outgoing strip system.

\begin{proposition}\label{prop:fold-character}
Under the analytic hypotheses,
\[
 \chi^{\bar{\widetilde p}}(S_h)=\chi^r(S_h).
\]
\end{proposition}
\begin{proof}
Choose a smooth reference configuration
$c_0$ on the doubled disk. Attach a
representative sphere on the incoming
strip end, then move its node through
the disk to the outgoing end. The two
resulting configurations lie in the
same relative class: the path of the
node and its trivialization explicitly
gives a homotopy. If a product target is
used in the folded description, perform
this operation on the appropriate
factor before unfolding. The other
factor is kept fixed.

There is a continuous orientation on
the determinant bundle along this path.
Its construction on the folded domain
uses the collapse transport of
Section~\ref{sec:collapse}, the
closed-cut conversion, and
Hypothesis~\ref{hyp:completion}. At the
two ends, coherence identifies it with
the strip-node gluing orientations.
No Koszul sign is introduced by moving
the complex node factor. Their ratios
are therefore equal, giving the
assertion.

This comparison uses arbitrary smooth
configurations, not only holomorphic
solutions. To justify coherence where
the fold is attached to prescribed
strip configurations, select one
reference configuration for each of the
finitely many generator pairs and their
required sphere classes. Choose the
height and neck lengths so that
ordinary gluing applies to this finite
set. The determinant-line construction
then extends over each configuration
component by path independence. The
regularity and completion hypotheses
are needed on all the domains in this
finite construction; a comparison only
at the constant solution would not
suffice.
\end{proof}

Changing the almost-complex base on the
closed cut changes the output generator
lines by a fixed diagonal conversion.
That conversion cancels in a sphere
ratio. Thus the character comparison is
unchanged by the conversion required in
Section~\ref{sec:excision}, although that
conversion remains essential for the
cup identity as a chain-level formula.

\begin{corollary}[Reflection invariance]\label{cor:reflection-character}
For every regular datum $p$ in the
admissible ball at an object with
nonzero integral gauge homology,
\[
 \chi^{\bar p}(S_h)=\chi^p(S_h).
\]
\end{corollary}
\begin{proof}
Apply Theorem~\ref{thm:character-constancy}
on $-M$ to compare $\bar p$ with
$\bar{\widetilde p}$. Apply
Proposition~\ref{prop:fold-character}
to compare the latter with $r$.
Apply character constancy on $M$ to
compare $r$ with $p$. The reflected
ball is the same ball under the
chosen reflection identification.
Gauge duality makes nonvanishing
for $M$ and $-M$ equivalent.
\end{proof}

\subsection{The reflection pairing}
Compare on the strips of $\bar p$ its
own transported system and the dual
system of~\eqref{eq:reflection-sign}.
Their ratio is constant on smooth
configuration components and
multiplicative under concatenation.
Choose one reference intersection
point. A configuration from it to
any other intersection determines
a sign $\kappa_p(x)$. Changing that
configuration changes its class by
a sphere; Corollary
\ref{cor:reflection-character} says
that the ratio on this sphere is
one. Hence the ratio on a strip
from $x$ to $y$ is
$\kappa_p(x)\kappa_p(y)$, independent
of the strip class.

\begin{proposition}\label{prop:reflection-pairing}
The map
\begin{equation}\label{eq:Phi-S}
 \Phi_p:\CF(\bar p)\longrightarrow\CF(p)^\vee,
 \qquad o_x^{\bar p}\longmapsto
                \kappa_p(x)(o_x^p)^\vee
\end{equation}
is an even chain isomorphism. The
configuration-level comparison
determines it up to one sign. The
pairings may be chosen so that
\begin{equation}\label{eq:Phi-symmetry}
 \Phi_{\bar p}=\Phi_p^\vee\ev.
\end{equation}
\end{proposition}
\begin{proof}
Equation~\eqref{eq:reflection-sign}
and the ratio just computed give
exactly the dual differential under
\eqref{eq:Phi-S}. The paired gauge
bases make the sum of the two
endpoint parities zero, so the map
is even. Any two comparisons have
constant ratio: configurations
exist between every pair of
intersection points, and the ratio
on all those configurations fixes
the relative generator signs.
Only the initial reference sign
is free.

Reflect twice on a strip of index
$n$ with endpoints $x,y$. The two
transpose signs multiply to
$(-1)^{n(|x|+|y|)}=(-1)^n$.
The corresponding change of the
paired basis is the sign in the
graded evaluation map. Thus the
pairing obtained by the second
reflection is exactly
$\Phi_p^\vee\ev$, after the common
reference sign is fixed. This
proves~\eqref{eq:Phi-symmetry} at
configuration level.
\end{proof}

A chain-map condition alone would
not determine $\kappa_p$ up to
one sign. If the differential
splits into two subcomplexes,
one can change the sign on one
summand and still obtain a chain
map. The preceding construction
uses all strip configurations,
including those which are not
solutions, to exclude that extra
freedom.

Put
\begin{equation}\label{eq:hat-Phi}
 \widehat\Phi_p=\Phi_p^\vee\ev:
   \CF(p)\longrightarrow\CF(\bar p)^\vee.
\end{equation}
This is the reflection pairing
which will be composed with the
cup map~\eqref{eq:check-cup}.

\section{Continuation and the diagonal compression cup}
\label{sec:continuation}

\subsection{Coherent orientations with prescribed ends}
A continuation datum from $p$ to $q$ has a strip-like
surface with Lagrangian boundary conditions varying
on a compact interval and equal to those of $p,q$
at its ends. The almost-complex and Hamiltonian
data are likewise constant outside that interval.
For generators $x,y$, let $\mathcal B_{p,q}(x,y)$
be its smooth configuration space, with the
prescribed exponential limits. A coherent
orientation system assigns relative orientations
to the determinant lines over these spaces and
commutes with attaching strips at either end.

The space of configurations in a fixed relative
homotopy class is path connected. One may first
choose a continuous relative homotopy, smooth it
relative to the fixed ends and labels, and then
make it constant beyond common end cutoffs. Local
charts are given by an exponential map for a
metric making the boundary labels totally
geodesic. The topology in
Section~\ref{sec:geometry} makes these paths
continuous. The relative classes are a torsor
for the same sphere group as fixed-label strips,
and the index changes by four under its generator.
In particular there is at most one index-zero
component for fixed endpoints.

\begin{lemma}[Rigidity of coherent systems]\label{lem:rigidity}
Two coherent continuation orientation systems
with the same end systems differ by one sign.
Such a system exists if and only if the sphere
characters at the ends agree.
\end{lemma}
\begin{proof}
The ratio of two systems is constant on each
configuration component. If a strip is attached
at either end, the common end orientation cancels,
so the ratio does not change. Choose one
reference configuration from $x_0$ to $y_0$.
For any $x,y$, choose strips from $x$ to $x_0$
and from $y_0$ to $y$. Attaching these strips
transports the ratio to a configuration with
ends $x,y$. Any difference in relative class
can be supplied by a sphere class on either
strip. Since the end systems are the same,
this also leaves the ratio unchanged. Thus
there is only one sign for all configurations.

For existence choose a reference orientation
and extend it by the same gluing procedure.
The determinant bundle on each configuration
component is orientable: the boundary
Lagrangians are the parallelizable compression-
body Lagrangians, and the remaining sphere
factors have complex orientation. The only
compatibility left when the gluing paths are
changed is moving a sphere class from the
incoming end to the outgoing end. The two
extensions agree precisely when the two sphere
characters agree. This also proves necessity
by moving the node of one sphere through a
continuation configuration. The construction
is continuous because the determinant line
is a line bundle over the smooth
configuration space.
\end{proof}

In an index-zero comparison a finite form of
the lemma suffices. Choose one strip for each
input generator and one for each output
generator, in the classes needed to reach
the index-zero component. Coherence on that
finite collection fixes the ratios of all
zero-dimensional counts. This form permits
the height and neck thresholds to be chosen
after a finite set of configurations, without
a claim of uniform estimates over the whole
configuration space.

Theorem~\ref{thm:character-constancy} supplies
existence for $p,q$ in the same admissible
ball. For a lower datum and an auxiliary upper
height datum, the folded configuration itself
supplies a coherent system and hence the
required character equality. The latter
observation applies even when the auxiliary
upper datum is outside the character-constancy
ball, provided it satisfies the separate
primary regularity threshold.

\subsection{Parity and continuation maps}
For a continuation configuration $c$ from $x$
to $y$, the number
\[
 d(p,q)=|y|-|x|-\ind D_c\pmod2
\]
is independent of the configuration and its
end generators. It is additive under
composition of continuation data. The
configuration description by sphere classes
shows its independence of the chosen datum.
For the paired gauge bases of this paper,
$d(p,q)=0$ at the data used in the comparison.
Here one can use the integral cup formula to
check the parity without assuming that
integral nonvanishing implies mod-two
nonvanishing.

Indeed, at a fold pair $(p,r)$ the map
$B(p,r)$ is homogeneous and
\eqref{eq:cup-zero} identifies it up to
homotopy with the even isomorphism
$N_r\cont_G(p,r)n_p$. On nonzero integral
homology this map cannot be both even and
odd. Thus its degree is zero. The doubled
configuration has the same index as the
continuation problem, and the gauge
reflection pairing is even; hence
$d(p,r)=0$. A common upper height datum
and additivity give $d(p,q)=0$. All uses
of the integral cup formula in this
argument precede the diagonal-cup
identification and therefore do not use
the continuation properties being proved.

With a coherent system define
\[
 \cont_S(p,q):\CF(p)\longrightarrow\CF(q)
\]
by the signed zero-dimensional continuation
count. It is even and commutes with the
differential. The compactification of a
one-dimensional continuation moduli space
proves this exactly as in
Proposition~\ref{prop:N-chain}.

\begin{proposition}\label{prop:continuation}
The continuation maps satisfy, up to sign
and chain homotopy,
\begin{align}
 \cont_S(q,r)\cont_S(p,q)&\simeq
                 \pm\cont_S(p,r),\label{eq:cont-composition}\\
 \cont_S(p,p)&\simeq\pm1.\label{eq:cont-unit}
\end{align}
Their homotopy classes are independent of
the continuation datum. They are homotopy
equivalences.
\end{proposition}
\begin{proof}
For a long strip neck, ordinary gluing
identifies the zero-dimensional solutions
with pairs of continuation solutions.
The determinant gluing is coherent, so
the signed matrix is the composite matrix.
The glued orientation system has the same
end systems as the chosen direct
continuation system. Lemma~\ref{lem:rigidity}
compares them by one sign. Varying the
neck length gives the chain homotopy in
\eqref{eq:cont-composition}.

A homotopy of continuation data, constant
near the ends, gives the corresponding
homotopy by its oriented one-dimensional
boundary. The pure symplectic compactness
and gluing used here are included in
Hypothesis~\ref{hyp:regularity}(d). For
constant data the zero-dimensional
solutions are constant strips, with
canonical positive orientation. This
normalizes the unit. Apply composition
to a datum and its reverse to obtain a
homotopy inverse.
\end{proof}

\subsection{Reflection of continuation}
There is no translation quotient for a
continuation configuration. Reflecting it
therefore has the graded transpose sign
but no additional translation sign.
Using the pairings of
Proposition~\ref{prop:reflection-pairing}
gives
\begin{equation}\label{eq:cont-reflection}
 \cont_S(\bar p,\bar q)\simeq
 \pm\Phi_q^{-1}\cont_S(q,p)^\vee\Phi_p.
\end{equation}
To verify coherence of the reflected
system, glue an input strip, a
continuation configuration, and an
output strip, and reflect the three
pieces in reverse order. If their
indices are $n_u,n_c,n_w$, the reordering
sign is
$(-1)^{n_un_c+n_un_w+n_cn_w}$.
Combining it with the three transpose
signs leaves $(-1)^{n_ud(p,q)}=1$.
The generator corrections $\kappa_p$
and $\kappa_q$ identify the reflected
end systems with the own transported
systems. Rigidity then gives the one
sign in~\eqref{eq:cont-reflection}.
This argument also shows why parity
zero is needed here.

\subsection{Identification of the diagonal cup}
For $k=0$ the compression correspondence
is the diagonal. Double the folded
triangle, glue the two distinguished
caps, and unfold along that diagonal.
The resulting surface is a continuation
strip. This is the geometric
identification used in Paper~I. To
obtain an integral statement we compare
its coherent systems before counting.

\begin{proposition}[Diagonal cup]\label{prop:diagonal-cup}
At every admissible fold pair $(p,r)$,
\begin{equation}\label{eq:diagonal-cup}
 B(p,r)\widehat\Phi_p\simeq
                  \pm\cont_S(p,r).
\end{equation}
\end{proposition}
\begin{proof}
The doubled triangle has the transported
fold orientation. Attach the two
auxiliary cap configurations, each
normalized by its unit relative
invariant. Unfolding identifies the
resulting real Cauchy--Riemann operator
with the continuation operator, with
its reflected input paired through
$\widehat\Phi_p$. The normal direction
at the diagonal and the ordered input
lines give precisely the graded
transpose convention used in
\eqref{eq:hat-Phi}.

On the finite configuration set from
Lemma~\ref{lem:rigidity}, ordinary
strip gluing and the orientation
comparisons of the collapse and
folded domains show coherence with
the prescribed input and output
systems. Attach input and output
strips successively; this reduces
product-strip coherence to the
one-factor comparisons. After
extending over the configuration
components, the pulled-back system
is a continuation system. Rigidity
compares it with the chosen system
by one sign.

For sufficiently long cap necks,
the gluing theorem gives a bijection
of zero-dimensional solutions and
preserves these orientations. A cap
of index zero has one regular
solution with unit coefficient, so
no extra multiplicity occurs. The
counts are therefore the same
up to that sign. Varying the cap
length and the triangle data to the
ones defining $B(p,r)$ gives a
chain homotopy through the
one-dimensional moduli spaces.
The preglue-to-solution path is
within the pure section chart,
and determinant transport on that
path preserves the canonical
orientation of a regular solution.
This is the pure symplectic gluing
input in Hypothesis~\ref{hyp:regularity}.
\end{proof}

This completes Theorem~\ref{thm:C}.
The statement uses the compression
cup of the manuscript. It does
not replace continuation by an
inverse of a chosen mixed map.

\subsection{Objects with vanishing homology}\label{subsec:zero}
Suppose $H(C_G(M))=0$. Lemma~\ref{lem:acyclic}
and Proposition~\ref{prop:N-isomorphism}
make both $C_G(M,p)$ and $\CF(M,p)$
contractible for every presentation.
Every closed map from or to a
contractible complex is a boundary:
compose it with a contracting
homotopy and use
\eqref{eq:compose-homotopy}. Hence
there is exactly one morphism
between it and any other object
in the homotopy category.

At such an object we take the
continuation and reflection maps
needed for the comparison to be
zero representatives of those
unique morphisms. In particular
a reflection map here is a
homotopy equivalence between zero
objects, not an invertible matrix.
We assert neither character
constancy nor existence of the
same configuration-level reflection
pairing at this object. Every
identity involving it is interpreted
in the homotopy category.

This convention also handles a
word passing through an acyclic
intermediate object: its class is
zero, as is the corresponding
gauge composite. Reduction modulo
two preserves a contracting
homotopy. Thus Paper~I's maps
incident to that object are
null-homotopic as well.

Nonzero mod-two gauge homology
implies nonzero integral gauge
homology, since an acyclic finite
free integral complex would be
contractible. The standard mod-two
framed calculation therefore
places rational homology spheres
and $\#^k(S^1\times S^2)$ in the
nonvanishing case
\cite[Corollary~1.4 and Section~7.8]{Scaduto}.
We do not use the converse
implication: a nonzero finite
odd-torsion integral group need
not be detected modulo two.

\section{Choices, primary naturality and handle maps}
\label{sec:choices}

\subsection{Order of the auxiliary choices}
All choices in a fixed comparison involve
finitely many presentations. The following
order avoids requiring a perturbation to
be simultaneously small relative to a
parameter chosen after it.
\begin{enumerate}[label=(\arabic*)]
\item Fix the splittings, metric data,
paired gauge bases and trace-coordinate
families for all objects and their
reflections. Choose a common primary
regularity threshold $r_H$ for this
finite set, as in
Proposition~\ref{prop:small-ball}.
\item Choose the first auxiliary heights
$t_*,t'_*$ so their height perturbations
and small generic corrections have norm
less than $r_H$. These data are used as
upper auxiliary presentations.
\item Choose the character-constancy
radius $r_*$ below $r_H$ and inside the
small-primary neighborhoods required
for those heights. Use the same choice
for a presentation and its reflection.
\item Choose smaller heights $t''$ for
the direct duality argument, $t_f$ for
the reflection fold, and $t_C$ for the
compression cup. Their height data
lie inside the corresponding balls.
Choose the lower datum of a cup after
$t_C$, sufficiently small relative to
that height.
\item Select the finite reference
configurations needed for coherence.
Choose their necks and the collapse
scales, then the regularizing data,
below all thresholds for those fixed
geometric data. In particular, an
auxiliary upper datum is a full
presentation with its own vertex-map
thresholds, not merely a pair of
transverse Lagrangians.
\item Choose the weights inside all
end gaps in that instance, below the
minimum of their decay rates. The
completed-domain estimates are required
at the same choices by
Hypothesis~\ref{hyp:regularity}.
\end{enumerate}
Every choice has a positive threshold
once the preceding finite data have
been fixed. Auxiliary upper data may
lie outside the final character ball;
the folded comparison supplies their
coherent continuation system directly.
The lower data and the handle target
where reflection is used lie inside
the relevant ball.

\subsection{Small-primary naturality}
We now prove the comparison for changes
of primary data. All maps in the next
two propositions are even.

\begin{proposition}\label{prop:primary-naturality}
For regular $p,q$ in one admissible ball,
\begin{equation}\label{eq:primary-naturality}
 N_q\cont_G(p,q)\simeq
              \pm\cont_S(p,q)N_p.
\end{equation}
\end{proposition}
\begin{proof}
Choose one upper height presentation $a$
of $-M$ generic for $\bar p$ and
$\bar q$. Apply the cup formula and
Proposition~\ref{prop:diagonal-cup} to
both fold pairs. They give
\begin{align*}
 \cont_S(\bar p,a)&\simeq
 \pm N_a\cont_G(\bar p,a)
                       n_{\bar p}\widehat\Phi_{\bar p},\\
 \cont_S(\bar q,a)&\simeq
 \pm N_a\cont_G(\bar q,a)
                       n_{\bar q}\widehat\Phi_{\bar q}.
\end{align*}
All factors being eliminated are
homotopy equivalences. Cancel the
common upper datum using continuation
composition. Thus
\begin{align*}
 \cont_S(\bar p,\bar q)\simeq\pm
 &(n_{\bar q}\widehat\Phi_{\bar q})^{-1}
       \cont_G(\bar p,\bar q)\\
 &\hspace{30mm}\cdot
       n_{\bar p}\widehat\Phi_{\bar p}.
\end{align*}
Use the reflection formula
\eqref{eq:cont-reflection} on the
symplectic side and its gauge analogue
from~\eqref{eq:gauge-duality}. The
symmetric conventions give
\[
 \widehat\Phi_{\bar p}=\Phi_p,
 \qquad \Phi_{G,M}n_{\bar p}=N_p^\vee,
\]
and the same equalities at $q$.
Cancelling these pairings transforms
the displayed identity into
\[
 \cont_S(q,p)^\vee\simeq
 \pm\bigl(N_p\cont_G(q,p)N_q^{-1}\bigr)^\vee.
\]
Transpose and rename $p,q$ to obtain
\eqref{eq:primary-naturality}.
Every double-dual map occurs together
with its inverse. Because the maps
here are even, the cancellations
introduce no grading-dependent sign.
\end{proof}

This proof uses the same pairing in
the diagonal-cup and reflection
identities. Choosing an unrelated
chain-level duality isomorphism would
leave a diagonal automorphism after
the cancellation.

\subsection{Mixed duality}
Put $N_p^\dagger=n_p\widehat\Phi_p$.
This is the transpose of the reflected
mixed map, with the gauge and
symplectic pairings inserted.

\begin{proposition}\label{prop:mixed-duality}
For every regular $p$ in the admissible
ball,
\begin{equation}\label{eq:mixed-duality}
 N_pN_p^\dagger\simeq\pm1.
\end{equation}
\end{proposition}
\begin{proof}
First take $p$ small enough for a fold
at the height $t''$, with upper datum
$r$ in the same ball. The two cup
identities give
\[
 \cont_S(p,r)\simeq
 \pm N_r\cont_G(p,r)n_p\widehat\Phi_p.
\]
By primary naturality the right side
is homotopic, up to sign, to
$\cont_S(p,r)N_pN_p^\dagger$.
Cancel the continuation equivalence.

For a general datum $q$ in the ball,
transpose primary naturality for the
reflected presentations. The symmetric
pairings give
\[
 N_q^\dagger\cont_S(p,q)\simeq
             \pm\cont_G(p,q)N_p^\dagger.
\]
Combine this with primary naturality
and~\eqref{eq:mixed-duality} at $p$:
\[
 N_qN_q^\dagger\cont_S(p,q)
             \simeq\pm\cont_S(p,q).
\]
Cancel the continuation once more.
The contractible case has the
interpretation of
Section~\ref{subsec:zero}.
\end{proof}

\subsection{One- and three-handle squares}
At the compatible handle data define
the integral compression cup by
\begin{equation}\label{eq:U}
 U_C=\check U_C\widehat\Phi_p:
               \CF(p)\longrightarrow\CF(q),
 \qquad q=C\circ p.
\end{equation}
Define the cap by reflected
transposition. If $\bar C$ is the
reflected compression and
$\check U_{\bar C}:\CF(p)^\vee\to\CF(\bar q)$
is its unpaired cup, put
\begin{equation}\label{eq:V}
 \check V_C=\ev_{\CF(p)}^{-1}
                  (\check U_{\bar C})^\vee,
 \qquad V_C=\check V_C\widehat\Phi_q.
\end{equation}
Then $V_C:\CF(q)\to\CF(p)$. Equivalently
it is the transpose of the paired
reflected cup with the two pairings
inserted. These are the orientations
of the compression cup and cap counts
of Paper~I. At other presentations
they are conjugated by the geometric
continuation maps.

\begin{theorem}\label{thm:handles}
For the one-handle cobordism $W_k$ and
its corresponding three-handle
cobordism $W_3$, one has
\begin{align}
 U_CN_p&\simeq\pm N_q\IG(W_k),
                                      \label{eq:onehandle}\\
 V_CN_q&\simeq\pm N_p\IG(W_3).
                                      \label{eq:threehandle}
\end{align}
The identities hold at all admissible
endpoint presentations.
\end{theorem}
\begin{proof}
The cup formula and mixed duality give
\[
 U_CN_p\simeq
 \pm N_q\IG(W_k)n_p\widehat\Phi_pN_p
 \simeq\pm N_q\IG(W_k),
\]
which proves~\eqref{eq:onehandle}.
For the cap, transpose the cup formula
at the reflected compression. The
graded evaluation maps and
$\Phi_{G,-M}=\Phi_{G,M}^\vee\ev$
give the gauge transpose of the
reversed handle map. Mixed duality
at $q$ cancels the reflected mixed
map. In detail, in the homotopy
category,
\begin{align*}
 V_CN_q
 &=\check V_C\widehat\Phi_qN_q\\
 &\simeq\pm\check V_C n_q^{-1}\\
 &\simeq\pm N_p\Phi_{G,-M}^{-1}
            \IG(\overline W_3)^\vee
            \Phi_{G,-M^+}\\
 &\simeq\pm N_p\IG(W_3).
\end{align*}
The only potentially odd map in the
transpose comparison is the handle
map; all pairings and mixed maps
are even. Its degree is retained
in the graded transpose, so any
remaining sign is scalar.

Finally conjugate both handle maps
by continuations from arbitrary
admissible endpoints to the
compatible handle data. Primary
naturality moves the mixed maps
past these continuations, and
gauge functoriality identifies the
gauge composite with the same
cobordism. Homotopies compose by
\eqref{eq:compose-homotopy}.
\end{proof}

\section{Descent to cobordisms and proof of the main theorem}
\label{sec:descent}

\subsection{The remaining changes of presentation}
The geometric presentation moves are the
ones in Paper~I. We record how the
integral orientations enter them.
A change of regularizing data uses
the family defining
Proposition~\ref{prop:N-isomorphism}.
A change of strip-end almost-complex
structure uses the analogous mixed
continuation family; the end systems
are identified by the determinant
transport in
Proposition~\ref{prop:transport}.
For a change of gauge base, both the
gauge and symplectic generator lines
are converted by the same
right-placed maps. The reference
constant pair has no extra line,
so its orientation comparison
commutes with this conversion.

An isomorphism of decorated data
pulls back the equation, its
linearization and the shifting
construction. Choose pulled-back
stabilizations in the determinant
construction. The kernel gluing
maps then commute with the
isomorphism, and the mixed square
commutes exactly on chains at the
pulled-back data. A subsequent
continuation gives the square at
the originally chosen data, up to
sign and homotopy. This applies
to frame changes and diffeomorphisms
preserving the torus decoration.

A Heegaard stabilization is realized
by the cancelling one- and two-handle
pair of Paper~I. Compose
\eqref{eq:onehandle} and
\eqref{eq:twohandle}. The decorated
gauge cobordism is the identity,
so its map is continuation up to
sign and homotopy. The reverse
stabilization is treated by the
cancelling two- and three-handle
pair. Passing through a common
stabilization treats changes of
metric data and isotopies of the
splitting. The order of choices
in Section~\ref{sec:choices} allows
the new vertices to satisfy the
integral height and ball conditions.
No orientation conclusion is
inferred from cancellation of
unsigned counts alone.

Thus every elementary presentation
arrow $e:p\to q$ has a symplectic
map $S_e$ constructed from the
curves of Paper~I and a square
\begin{equation}\label{eq:edge-square}
 N_qG_e\simeq\pm S_eN_p.
\end{equation}
The gauge arrows $G_e$ represent
the relevant decorated cobordism
or identity continuation. This
includes both even presentation
moves and possibly odd handle
arrows.

\subsection{Projective descent}
\begin{lemma}\label{lem:descent}
Suppose a connected presentation
category maps to a cobordism
category and every cobordism has
an elementary word. Let $G_e$
be a projective gauge functor on
these words. Suppose $N_p$ are
isomorphisms in a homotopy
category and the elementary
maps satisfy
\eqref{eq:edge-square}. Then
$S_d$ depends only on the
represented cobordism and the
endpoint presentations, up to
sign and homotopy.
\end{lemma}
\begin{proof}
For a word $d=e_m\cdots e_1$,
compose the squares to obtain
\[
 S_dN_{p_0}\simeq
     \pm N_{p_m}G_{e_m}\cdots G_{e_1}.
\]
Every composition of homotopies
is valid by
\eqref{eq:compose-homotopy}.
If $d'$ represents the same
cobordism, gauge functoriality
identifies the two gauge words
up to sign and homotopy. Cancel
$N_{p_0}$ to identify $S_d$
and $S_{d'}$.

For words representing the
identity cobordism this also
gives the transitive system of
comparison maps between
presentations. The identity
word is the identity map, and
concatenation of words gives
composition up to sign. Hence
these systems descend to a
projective functor and the
classes of $N_p$ give a
projective natural isomorphism.
\end{proof}

\begin{proof}[Proof of Theorem~\ref{thm:A}]
Proposition~\ref{prop:N-isomorphism}
gives the integral mixed
isomorphisms. Equation
\eqref{eq:twohandle},
Theorem~\ref{thm:handles}, and
the presentation comparisons
above supply all elementary
squares. Generation and
connectivity are the geometric
presentation results of
Paper~I, with the additional
smallness conditions arranged
as in Section~\ref{sec:choices}.
Proposition~\ref{prop:gauge-functor}
and Lemma~\ref{lem:descent}
therefore give
\eqref{eq:main} and the
projective functoriality.

The mixed maps preserve the
chosen object gradings. On
homology, the comparison with
the homogeneous gauge map
therefore gives the asserted
parity of the symplectic map.
No separate assertion of
homogeneity for every chosen
chain representative of a
long word is needed.

Modulo two all determinant
orientations give the same
counts. The conversions become
the identity on the generator
basis, the signs in the tensor
and dual formulas disappear,
and the pairing
\eqref{eq:Phi-S} becomes the
paired-basis reflection
pairing. Thus the geometric
maps reduce to those of
Paper~I at the same data. At
an acyclic object the zero
representatives reduce to
maps homotopic to those of
Paper~I by
Section~\ref{subsec:zero}.
\end{proof}

\subsection{What remains for an unconditional extension}
Theorem~\ref{thm:A} is a
conditional integral comparison.
Removing its analytic hypotheses
requires a proof on the actual
completed mixed domains of
uniform regularity, uniform
slice and weight estimates,
and continuity of the
orientation identifications
in the torus completion.
The nonseparating self-gluing
statement must also be verified
with the lifting groups and
ordered determinant lines used
here. These steps are separate
from the orbit-count algebra in
Proposition~\ref{prop:capped}.

Several further extensions ask
different questions. Exact
composition signs require more
than projectivizing the maps;
the regrading discussed after
\eqref{eq:gauge-dictionary}
indicates the extra grading
information involved. An
identification with integral
homology-oriented instanton
maps requires a comparison of
orientation systems on the
framed quotient. An
identification with integral
quilted Floer maps requires a
comparison with relative-spin
orientations on the compression
correspondences. None of these
identifications follows from
the conditional theorem by
reducing modulo two.

\appendix
\section{Analytic inputs and their precise uses}\label{app:inputs}

This appendix records the analytic
statements used in the draft and
separates the ordinary inputs from
the completed-domain hypotheses.
All references to Paper~I mean
its version of 25 September 2026.
The present paper does not
supply an independent proof of
that manuscript's unsigned
collapse, torus excision, or
generation results.

\subsection{Mixed Fredholm theory and shifting}
The mixed linearization couples
the anti-self-dual operator to a
real Cauchy--Riemann operator by
the linearized matching condition.
For fixed admissible ends it is
Fredholm on the prescribed
weighted spaces. The local
regularity and compactness
statements for the original
mixed equation are in
\cite{DFL-mixed}; the vertex
Fredholm and constant-pair
statements used here are in
\cite[Theorems~4.12, 4.17 and
Proposition~3.19]{DFL}.
The mixed shifting statements
are \cite[Lemmas~4.62 and~4.64]{DFL}.
Their use on the completed and
folded domains, including the
relative family form in the
smooth topology, is explicitly
covered by
Hypothesis~\ref{hyp:regularity}.

At an ordinary strip end the
asymptotic operator is
$J\partial_t$ with the totally
real endpoint conditions. At a
gauge end it is the extended
Hessian. A weight $\delta$ must
be separated from the relevant
spectrum; at a decaying end it
must be smaller than the decay
rate of the configuration and
of the linearized kernel and
cokernel elements. At a mixed
end the stationary modes are
split off before the weight is
imposed. This last finite
factor remains in determinant
gluing and may not be discarded.

\subsection{Boundary Hodge theory and holonomy regularity}
For a flat orthogonal bundle on
a compact compression body,
the input needed in
Proposition~\ref{prop:small-ball}
is the Hodge decomposition on
$W^{1,2}$ with normal boundary
condition. Its harmonic fields
are smooth and finite
dimensional, and its exact
part is the covariant
differential of $W^{2,2}$
sections. The Dirichlet Green
operator for $d_B^*d_B$
preserves smoothness, so
$1-d_BG_Dd_B^*$ is a
smoothness-preserving
projection onto coclosed forms.
These statements follow from
the flat covariant de Rham
case of
\cite[Section~1.6 and
Theorems~3.10--3.11]{KL}; the
normal and tangential versions
are interchanged by Hodge star.

The boundary elliptic estimate
has the form
\[
 \|b\|_{W^{1,p}}\le C
 (\|d_Bb\|_{L^p}+\|d_B^*b\|_{L^p}
                              +\|b\|_{L^p})
\]
for $*b|_{\partial C}=0$.
For smooth fixed $B$, adjoint
bundle coefficients add only
bounded zeroth-order terms in
local trivializations. The
compactness and Coulomb-slice
inputs are Theorems A, D(a)
and F of the introductory
statements in
\cite{Wehrheim-book}.

Finite holonomy functions have
bounded gradients and
Lipschitz gradient differences
as used in
\eqref{eq:primary-norm}
\cite[Proposition~6.1]{DFL}.
Their higher derivatives obey
the local tame estimates
obtained by differentiating
parallel transport. Combined
with the boundary elliptic
estimate, these give the
bootstrap from the slice
fixed point to the smooth
perturbed-flat solution.
For holonomy families with
moving basepoints, restriction
and composition are performed
in the smooth topology, not
asserted bounded on arbitrary
low-order Sobolev spaces.

\subsection{Pure symplectic gluing}
The continuation and cap
arguments require the
following form of gluing.
For a finite collection of
regular configurations with
nondegenerate strip ends,
there is a sufficiently long
neck such that every nearby
zero-dimensional solution is
the unique correction of a
preglued configuration. The
correction tends to zero with
the neck length, and the
linearized gluing map
identifies the determinant
orientations. The estimates
are uniform over compact
families of the fixed
regular configurations.
The strip construction is
provided by the gluing
analysis of~\cite{Simcevic};
for quilted surfaces see
\cite{Mau,MWW}. Exponential
decay is used at every
ordinary end; in the
transverse case the full-strip
statement is
\cite[Proposition~2.3.1]{Simcevic}.

Compactness in dimensions
zero and one requires the
monotonicity and index
restrictions for the actual
Lagrangians and
correspondences. The
compactified one-dimensional
space then has the endpoint
and ordinary broken-strip
strata used for continuation.
The monotone quilted
compactness and family
constructions are in
\cite[Theorems~3.9 and~4.1]{WW-relative}
and~\cite[Section~3]{MWW}.
Their applicability to every
pure symplectic domain used
here is included in
Hypothesis~\ref{hyp:regularity}(d).
A cohomological composition
identity alone is not used
as a substitute for this
moduli-space gluing statement.

\subsection{The unresolved uniform comparisons}
The following are the forms
of the remaining estimates
needed to remove the
hypotheses.
For the completed gauge
slices one needs a uniform
Poincare estimate
\[
 \|\zeta\|_{L^2}\le C\|d_A\zeta\|_{L^2}
\]
on the specified gauge
subspace, including the
end framing conditions.
Together with the higher
elliptic estimates it must
give a positive Coulomb
slice radius independent of
the growing torus length.
For mixed linearizations
one needs uniform
regularity on the joined,
separate and closed-cut
domains, with the same
asymptotic weights and with
the finite-dimensional
stationary coordinates
retained. Finally the
preglued-to-solution paths
must lie in the
configuration charts where
the determinant transport
is defined. Pointwise
regularity at each finite
length proves none of
these uniform assertions
by itself.

For the gauge self-gluing,
one must additionally check
that the compactified
one-dimensional stretching
family has no omitted
end, and that the kernel
of the restriction of
lifting classes is exactly
the gluing fiber in
Lemma~\ref{lem:gluing-multiplicity}.
The finite group
calculation in
Section~\ref{sec:excision}
then gives integral unit
multiplicity. Keeping these
analytic statements
separate makes clear what
is assumed by the first
draft and what the
orientation algebra proves
from them.

\section{A sign calculation for projective composition}\label{app:signs}

The conversion of a strictly
commuting map $F$ of degree
$f$ to a closed map is
$F^{\mathrm r}=F\eps^f$.
For another strictly
commuting map $G$ of degree
$g$, the relation
$\eps^fG=(-1)^{fg}G\eps^f$
gives
\[
 F^{\mathrm r}G^{\mathrm r}
   =(-1)^{fg}FG\eps^{f+g}
   =(-1)^{fg}(FG)^{\mathrm r}.
\]
Thus this convention change
is compatible with
composition after
projectivization. The
factor is scalar, despite
the grading operator used
in converting each map.

For a strict homotopy
$F-F'=dH+Hd$, where
$|H|=f+1$, the correctly
converted homotopy is
$H\eps^f$. Indeed
\[
 \delta(H\eps^f)
 =dH\eps^f-(-1)^{f+1}H\eps^fd
 =(dH+Hd)\eps^f.
\]
Using $H\eps^{f+1}$ instead
would give the wrong sign.
This calculation is useful
whenever a gauge
composition homotopy is
translated into the Hom
complex convention.

Lastly, if
$F_i-G_i=\delta K_i$ and
the maps are closed,
a homotopy between two
composites is obtained
by replacing one factor
at a time and using
\eqref{eq:compose-homotopy}.
For example, for three
factors one may take
\[
 K_3F_2F_1
 +(-1)^{|G_3|}G_3K_2F_1
 +(-1)^{|G_3|+|G_2|}G_3G_2K_1.
\]
Applying $\delta$ telescopes
to $F_3F_2F_1-G_3G_2G_1$.
This is the chain-level
algebra used in
Lemma~\ref{lem:descent}.

\begin{thebibliography}{99}

\bibitem{DFL}
A.~Daemi, K.~Fukaya and M.~Lipyanskiy,
\emph{Lagrangians, $SO(3)$-instantons and the Atiyah--Floer conjecture},
arXiv:2109.07032.

\bibitem{DFL-mixed}
A.~Daemi, K.~Fukaya and M.~Lipyanskiy,
\emph{Lagrangians, $SO(3)$-instantons and mixed equation},
Geom. Funct. Anal. \textbf{34} (2024), 659--732;
arXiv:2109.07038.

\bibitem{Donaldson87}
S.~K. Donaldson, \emph{The orientation of Yang--Mills moduli spaces
and $4$-manifold topology}, J. Differential Geom. \textbf{26}
(1987), 397--428.

\bibitem{Donaldson-book}
S.~K. Donaldson, \emph{Floer Homology Groups in Yang--Mills Theory},
Cambridge Tracts in Math. \textbf{147}, Cambridge Univ. Press, 2002.

\bibitem{Floer-Hofer}
A.~Floer and H.~Hofer, \emph{Coherent orientations for periodic
orbit problems in symplectic geometry}, Math. Z. \textbf{212}
(1993), 13--38.

\bibitem{KM-unknot}
P.~B. Kronheimer and T.~S. Mrowka,
\emph{Khovanov homology is an unknot-detector},
Publ. Math. Inst. Hautes \'Etudes Sci. \textbf{113}
(2011), 97--208.

\bibitem{KM-sutures}
P.~B. Kronheimer and T.~S. Mrowka,
\emph{Knots, sutures, and excision},
J. Differential Geom. \textbf{84} (2010), 301--364.

\bibitem{KL}
R.~Kupferman and R.~Leder,
\emph{Elliptic pre-complexes, Hodge-like decompositions and
 overdetermined boundary-value problems}, arXiv:2304.08977v3.

\bibitem{Mau}
S.~Ma'u, \emph{Gluing pseudoholomorphic quilted disks},
arXiv:0909.3339.

\bibitem{MWW}
S.~Ma'u, K.~Wehrheim and C.~T. Woodward,
\emph{$A_\infty$ functors for Lagrangian correspondences},
Selecta Math. (N.S.) \textbf{24} (2018), 1913--2002;
arXiv:1601.04919v5.

\bibitem{Scaduto}
C.~W. Scaduto, \emph{Instantons and odd Khovanov homology},
J. Topol. \textbf{8} (2015), 744--810.

\bibitem{Simcevic}
T.~Sim\v{c}evi\'c, \emph{A Hardy space approach to Lagrangian Floer
gluing}, Ph.D. thesis, ETH Z\"urich, Diss. ETH No.~22118, 2014;
arXiv:1410.5998.

\bibitem{Wehrheim-book}
K.~Wehrheim, \emph{Uhlenbeck Compactness}, EMS Series of Lectures in
Mathematics, European Mathematical Society, Z\"urich, 2004.

\bibitem{WW-orientations}
K.~Wehrheim and C.~T. Woodward,
\emph{Orientations for pseudoholomorphic quilts},
arXiv:1503.07803.

\bibitem{WW-relative}
K.~Wehrheim and C.~T. Woodward,
\emph{Pseudoholomorphic quilts},
J. Symplectic Geom. \textbf{13} (2015), 849--904.

\bibitem{WW-quilt}
K.~Wehrheim and C.~T. Woodward,
\emph{Quilted Floer cohomology},
Geom. Topol. \textbf{14} (2010), 833--902.

\bibitem{PaperI}
B.~J. Wuebben, \emph{Cobordism maps and Atiyah--Floer isomorphisms
for framed instanton homology~I}, manuscript, 25 September 2026,
\url{https://github.com/bwuebben/donaldson-atiyah-floer/tree/main/af-cob-1}.

\bibitem{Zinger}
A.~Zinger, \emph{The determinant line bundle for Fredholm operators:
construction, properties, and classification},
Math. Scand. \textbf{118} (2016), 203--268.

\end{thebibliography}
\end{document}
